\title{Επισκόπηση Λειτουργίας Συστήματος Μεταφοράς Ηλεκτρικής Ενέργειας}

\section{Εισαγωγή}
\label{kef-01-theoria}

Το Σύστημα Μεταφοράς Ηλεκτρικής Ενέργειας (ΣΜΗΕ) αποτελεί τη ραχοκοκαλιά ενός Συστήματος
Ηλεκτρικής Ενέργειας (ΣΗΕ). Η ασφαλής λειτουργία του απαιτεί τη συνεχή εξισορρόπηση μεταξύ
παραγωγής και ζήτησης και, παράλληλα, την τήρηση των προβλεπόμενων ορίων λειτουργίας, τόσο στην τρέχουσα κατάσταση όσο
και μετά από απρόβλεπτα συμβάντα. Το κεφάλαιο αυτό παρουσιάζει συνοπτικά το πώς οργανώνεται η
λειτουργία του ΣΜΗΕ, από τον προγραμματισμό των επόμενων ημερών έως την εποπτεία και τον
έλεγχο σε πραγματικό χρόνο.

Αρχικά περιγράφονται η δομή του ΣΜΗΕ, οι απαιτήσεις λειτουργίας καθώς και τα
διαδοχικά στάδια προγραμματισμού μέσω των αγορών ηλεκτρικής ενέργειας. Ακολουθούν οι καταστάσεις λειτουργίας, το κριτήριο N-1 και τα συστήματα εποπτείας και ελέγχου σε πραγματικό χρόνο. Έμφαση δίνεται στην εκτίμηση κατάστασης, η οποία υπολογίζει μέσω των διαθέσιμων μετρήσεων τις τάσεις και τις γωνίες όλων των ζυγών. Το κεφάλαιο ολοκληρώνεται με τα ευέλικτα συστήματα μεταφοράς
εναλλασσόμενου ρεύματος (Flexible AC Transmission Systems, FACTS) για τον έλεγχο των ροών
ισχύος και των τάσεων. Στα Παραδείγματα που ακολουθούν, η εκτίμηση κατάστασης εφαρμόζεται σε
μικρά δίκτυα με κώδικα που εκτελείται απευθείας στην ηλεκτρονική έκδοση.

\section{Ρόλος και δομή του Συστήματος Μεταφοράς}

Το ΣΜΗΕ είναι το δίκτυο υπερυψηλής και υψηλής τάσης που διασυνδέει τους σταθμούς παραγωγής με τα κέντρα κατανάλωσης και με τα γειτονικά
συστήματα. Το Ελληνικό Σύστημα Μεταφοράς Ηλεκτρικής Ενέργειας (ΕΣΜΗΕ) λειτουργεί στα επίπεδα
τάσης 400 kV και 150 kV. Διαχειριστής του είναι ο Ανεξάρτητος Διαχειριστής Μεταφοράς
Ηλεκτρικής Ενέργειας (ΑΔΜΗΕ), μέλος του Ευρωπαϊκού Δικτύου Διαχειριστών Συστημάτων Μεταφοράς
Ηλεκτρικής Ενέργειας (European Network of Transmission System
Operators for Electricity, ENTSO-E). Το σύστημα ανήκει στη
συγχρονισμένη περιοχή της Ηπειρωτικής Ευρώπης (Continental Europe, CE).

Το διασυνδεδεμένο σύστημα μεταφοράς των μελών του ENTSO-E απεικονίζεται στον διαδραστικό
χάρτη του Σχήματος~\ref{fig:entsoe-map}, όπου φαίνονται οι γραμμές μεταφοράς ανά επίπεδο
τάσης και οι διασυνδέσεις μεταξύ των συστημάτων.

\input{embeds/entsoe-map.md}

\subsection{Απαιτήσεις λειτουργίας}

Ο Διαχειριστής Συστήματος Μεταφοράς (ΔΣΜ· Transmission System Operator, TSO) οφείλει να
ικανοποιεί ταυτόχρονα τις ακόλουθες απαιτήσεις:

\begin{itemize}
\item \textbf{Ισοζύγιο ισχύος.} Η παραγωγή πρέπει να ισούται συνεχώς με τη ζήτηση, συμπεριλαμβανομένων των ωμικών απωλειών του δικτύου. Κάθε ανισορροπία καλύπτεται αρχικά από την κινητική ενέργεια των στρεφόμενων μαζών και εκδηλώνεται ως μεταβολή της συχνότητας. Για παράδειγμα, ένα
      έλλειμμα παραγωγής προκαλεί επιβράδυνση των σύγχρονων μηχανών και πτώση της συχνότητας.
\item \textbf{Επιχειρησιακή ασφάλεια.} Το σύστημα πρέπει να παραμένει εντός των προκαθορισμένων ορίων (π.χ. θερμικά όρια του εξοπλισμού, όρια τάσης, όρια ευστάθειας), τόσο στην τρέχουσα κατάσταση όσο και μετά από οποιοδήποτε μεμονωμένο απρόβλεπτο συμβάν (contingency), όπως η απώλεια γραμμής ή μονάδας. Η απαίτηση αυτή διατυπώνεται ως κριτήριο N-1, το οποίο ορίζεται στην υποενότητα
      «\nameref{sec:n1}».
\item \textbf{Ποιότητα ηλεκτρικής ενέργειας.} Τα ηλεκτρικά χαρακτηριστικά της τάσης (συχνότητα, πλάτος, αρμονικό περιεχόμενο, κτλ.) πρέπει να
      διατηρούνται εντός προδιαγεγραμμένων ορίων.
\end{itemize}

Οι παραπάνω απαιτήσεις πρέπει να ικανοποιούνται σε όλες τις χρονικές κλίμακες λειτουργίας ενός
ΣΜΗΕ. Για τον σκοπό αυτό υιοθετείται ένας σύνθετος και πολύπλευρος μηχανισμός ελέγχου, ο οποίος
αντιμετωπίζει τα φαινόμενα που εμφανίζονται σε κάθε κλίμακα, όπως συνοψίζει ο
Πίνακας~\ref{tab:chronikes-klimakes}. Στο πλαίσιο του μαθήματος εξετάζονται κυρίως οι χρονικές
κλίμακες από τα δευτερόλεπτα και άνω.

\begin{table}[!htbp]
\centering
\begin{tabular}{lll}
\hline
Χρονική κλίμακα & Φαινόμενο ή διαδικασία & Μηχανισμός \\
\hline
μs -- ms   & βραχυκυκλώματα, υπερτάσεις   & προστασία, διακόπτες ισχύος \\
ms -- s    & ηλεκτρομηχανικές ταλαντώσεις & αδράνεια, PSS, FACTS \\
s -- 30 s  & απόκλιση συχνότητας          & εφεδρεία διατήρησης (ΕΔΣ) \\
30 s -- 15 min & αποκατάσταση συχνότητας  & εφεδρεία αποκατάστασης (ΕΑΣ) \\
min -- h   & αναπρογραμματισμός           & ενδοημερήσια αγορά \\
h -- ημέρα & ένταξη μονάδων, κατανομή     & αγορά επόμενης ημέρας \\
εβδομάδες -- έτη & επάρκεια ισχύος, συντηρήσεις & προγραμματισμός, επενδύσεις \\
\hline
\end{tabular}
\caption{Χρονικές κλίμακες λειτουργίας του ΣΜΗΕ και αντίστοιχοι μηχανισμοί ελέγχου. Οι
σταθεροποιητές συστήματος ισχύος (Power System Stabilizers, PSS) αποσβένουν τις
ηλεκτρομηχανικές ταλαντώσεις μέσω του συστήματος διέγερσης των γεννητριών.}
\label{tab:chronikes-klimakes}
\end{table}

\section{Στάδια λειτουργίας: από τον προγραμματισμό στον πραγματικό χρόνο}
\label{sec:stages}

Ο προγραμματισμός της λειτουργίας γίνεται σε διαδοχικούς χρονικούς ορίζοντες. Καθώς
πλησιάζει ο χρόνος φυσικής παράδοσης, η αβεβαιότητα των προβλέψεων φορτίου και ανανεώσιμης
παραγωγής μειώνεται, ενώ οι διαθέσιμες ενέργειες διόρθωσης περιορίζονται και το κόστος τους
αυξάνεται. Κάθε στάδιο διορθώνει τις αποκλίσεις του προηγούμενου. Στο ευρωπαϊκό
Μοντέλο-Στόχο (Target Model) οι ορίζοντες αυτοί αντιστοιχούν σε διακριτές αγορές:

\begin{itemize}
\item \textbf{Μακροπρόθεσμος και προθεσμιακός ορίζοντας} (έτη έως εβδομάδες). Μελέτες
      επάρκειας ισχύος, συντονισμός των προγραμμάτων συντήρησης γραμμών και μονάδων,
      διαστασιολόγηση των εφεδρειών, συναλλαγές στην Προθεσμιακή Αγορά για αντιστάθμιση
      του κινδύνου τιμής.
\item \textbf{Αγορά Επόμενης Ημέρας} (day-ahead). Αποτελεί το κύριο στάδιο
      προγραμματισμού. Οι συμμετέχοντες υποβάλλουν εντολές αγοράς και πώλησης για κάθε
      αγοραία χρονική μονάδα της επόμενης ημέρας. Η αγορά εκκαθαρίζεται μέσω της ενιαίας
      σύζευξης επόμενης ημέρας σε ευρωπαϊκό επίπεδο (Single Day-Ahead Coupling, SDAC), με
      μεγιστοποίηση του συνολικού κοινωνικού πλεονάσματος υπό τους περιορισμούς της
      διαθέσιμης διαζωνικής δυναμικότητας· οι ροές μεταξύ ζωνών προσφοράς κατευθύνονται
      από ζώνες χαμηλότερης προς ζώνες υψηλότερης τιμής. Ο ΔΣΜ ελέγχει στη συνέχεια την
      τεχνική εφικτότητα του προγράμματος.
\item \textbf{Ενδοημερήσια Αγορά} (intra-day). Επιτρέπει στους συμμετέχοντες να
      αναπροσαρμόζουν τις θέσεις τους καθώς βελτιώνονται οι προβλέψεις. Η σημασία της
      αυξάνεται με τη διείσδυση αιολικής και φωτοβολταϊκής παραγωγής, επειδή το σφάλμα
      πρόβλεψης της παραγωγής αυτής μειώνεται σημαντικά με τη μείωση του ορίζοντα
      πρόβλεψης.
\item \textbf{Πραγματικός χρόνος και εξισορρόπηση} (balancing). Ο ΔΣΜ ενεργοποιεί προσφορές
      ενέργειας εξισορρόπησης για την κάλυψη των εναπομενουσών ανισορροπιών, την
      αντιμετώπιση διαταραχών και τη διαχείριση συμφορήσεων.
\end{itemize}

\subsection{Αγοραία χρονική μονάδα}
\label{sec:mtu}

Η \emph{αγοραία χρονική μονάδα} (Market Time Unit, MTU) είναι η χρονική ανάλυση εκκαθάρισης
της αγοράς. Από την ημέρα παράδοσης 1η Οκτωβρίου 2025 (δημοπρασία της 30ής Σεπτεμβρίου
2025), η ενιαία σύζευξη επόμενης ημέρας λειτουργεί με αγοραία χρονική μονάδα \textbf{15
λεπτών} αντί της ωριαίας, σε εναρμόνιση με την περίοδο εκκαθάρισης ανισορροπιών. Με
μικρότερη αγοραία χρονική μονάδα, οι ταχείες μεταβολές φορτίου και παραγωγής, όπως οι
κλίσεις της φωτοβολταϊκής παραγωγής, αποτυπώνονται στο πρόγραμμα αντί να καλύπτονται με
ενέργεια εξισορρόπησης· μειώνονται έτσι οι αποκλίσεις στις αλλαγές ωριαίων περιόδων και οι
ανάγκες σε εφεδρείες.

Σε κάθε αγοραία χρονική μονάδα οι εντολές πώλησης ταξινομούνται κατά αύξουσα τιμή και οι
εντολές αγοράς κατά φθίνουσα. Η τομή των δύο αθροιστικών καμπυλών ορίζει την \textbf{Τιμή
Εκκαθάρισης Αγοράς} (ΤΕΑ· Market Clearing Price, MCP) και την εκκαθαριζόμενη ποσότητα
(Σχήμα~\ref{fig:dam}α). Όλες οι εντολές που γίνονται αποδεκτές εκκαθαρίζονται στην ίδια τιμή
(οριακή τιμολόγηση, pay-as-clear): οι πωλητές με χαμηλότερη τιμή προσφοράς, όπως οι
ανανεώσιμες πηγές με μηδενικό μεταβλητό κόστος, λαμβάνουν την ΤΕΑ και όχι την τιμή της
προσφοράς τους. Στη συζευγμένη αγορά οι ΤΕΑ όλων των ζωνών προσφοράς υπολογίζονται ταυτόχρονα
από κοινό αλγόριθμο (Pan-European Hybrid Electricity Market Integration Algorithm,
EUPHEMIA) και διαφέρουν μεταξύ δύο ζωνών μόνο όταν η διαζωνική
δυναμικότητα που τις συνδέει εξαντλείται.

Η ημερήσια διακύμανση της ΤΕΑ ακολουθεί τη διαφορά ζήτησης και ανανεώσιμης παραγωγής: σε
ηλιόλουστες ημέρες η τιμή μειώνεται τις μεσημεριανές ώρες, όταν η φωτοβολταϊκή παραγωγή είναι
μέγιστη, και αυξάνεται απότομα το απόγευμα, όταν η παραγωγή αυτή μηδενίζεται και τη ζήτηση
καλύπτουν θερμικές και υδροηλεκτρικές μονάδες (Σχήμα~\ref{fig:dam}β). Με αγοραία χρονική μονάδα
15 λεπτών οι κλίσεις αυτές αποτυπώνονται στις τιμές και στα προγράμματα, ενώ η ωριαία μονάδα
τις εξομαλύνει. Οι πραγματικές τιμές της ελληνικής ζώνης προσφοράς δημοσιεύονται καθημερινά
από το ΕΧΕ, σε πίνακες αποτελεσμάτων και σε διαδραστικά
\href{https://www.enexgroup.gr/el/web/guest/day-ahead-market-figures}{διαγράμματα της Αγοράς
Επόμενης Ημέρας}· το διάγραμμα της ΤΕΑ είναι ενσωματωμένο στο Σχήμα~\ref{fig:henex-tea}.

\begin{figure}[!htbp]
\centering
\includegraphics[width=\linewidth]{figures/dam-clearing.svg}
\caption{Σχηματική απεικόνιση της Αγοράς Επόμενης Ημέρας. (α) Αθροιστικές καμπύλες εντολών
πώλησης και αγοράς σε μία αγοραία χρονική μονάδα· η τομή τους ορίζει την Τιμή Εκκαθάρισης
Αγοράς (ΤΕΑ) και την εκκαθαριζόμενη ποσότητα $Q^*$. (β) Ενδεικτική ημερήσια διακύμανση της
ΤΕΑ με αγοραία χρονική μονάδα 15 λεπτών και ωριαία. Οι τιμές είναι συνθετικές· οι πραγματικές
δημοσιεύονται από το ΕΧΕ.}
\label{fig:dam}
\end{figure}

\input{embeds/henex-tea.md}

\subsection{Ελληνική αγορά ηλεκτρικής ενέργειας}
\label{sec:greek-market}

Στην Ελλάδα το Μοντέλο-Στόχος εφαρμόζεται από την 1η Νοεμβρίου 2020. Οι αγορές οργανώνονται
σε δύο διαχειριστές:

\begin{itemize}
\item Το \textbf{Ελληνικό Χρηματιστήριο Ενέργειας} (ΕΧΕ) διαχειρίζεται την Αγορά Επόμενης
      Ημέρας και την Ενδοημερήσια Αγορά. Η Αγορά Επόμενης Ημέρας είναι συζευγμένη μέσω της
      SDAC με τις γειτονικές ζώνες προσφοράς, με την Ιταλία από τον Δεκέμβριο του 2020 και με
      τη Βουλγαρία από τον Μάιο του 2021, και δίνει για κάθε αγοραία χρονική μονάδα την ΤΕΑ της
      ελληνικής ζώνης προσφοράς. Το ενεργειακό μείγμα που προκύπτει από την αγορά για κάθε
      15λεπτη αγοραία χρονική μονάδα δημοσιεύεται στην αρχική σελίδα του ΕΧΕ
      (Σχήμα~\ref{fig:henex-mix}).
\item Ο \textbf{ΑΔΜΗΕ} διαχειρίζεται την \textbf{Αγορά Εξισορρόπησης}, η οποία εξασφαλίζει ότι
      τα προγράμματα της αγοράς είναι τεχνικά εφικτά και ότι το σύστημα διαθέτει επαρκείς
      εφεδρείες, και καλύπτει σε πραγματικό χρόνο τις εναπομένουσες ανισορροπίες.
\end{itemize}

\input{embeds/henex-mix.md}

Η Αγορά Εξισορρόπησης αποτελείται από τρεις διαδικασίες:

\begin{enumerate}
\item \textbf{Αγορά Ισχύος Εξισορρόπησης}, η οποία εκτελείται μέσω της \emph{Διαδικασίας
      Ενοποιημένου Προγραμματισμού} (ΔΕΠ· Integrated Scheduling Process, ISP). Η ΔΕΠ
      διατυπώνεται ως πρόβλημα μεικτού ακέραιου προγραμματισμού, στο οποίο ενέργεια και ισχύς
      εξισορρόπησης βελτιστοποιούνται από κοινού, με ελαχιστοποίηση του συνολικού κόστους υπό
      τους τεχνικούς περιορισμούς των μονάδων και του ΕΣΜΗΕ. Δεσμεύει έξι προϊόντα ισχύος
      εξισορρόπησης, ανοδική και καθοδική ΕΔΣ, αΕΑΣ και χΕΑΣ, δηλαδή ΕΑΣ με αυτόματη και με
      χειροκίνητη ενεργοποίηση (Κεφάλαιο 3), και καθορίζει τα προγράμματα κατανομής.
      Εκτελείται τρεις φορές ημερησίως, στις 16:45 της προηγούμενης ημέρας (ΔΕΠ1)
      και στις 00:00 (ΔΕΠ2) και 12:00 (ΔΕΠ3) της ημέρας παράδοσης, ώρα Ελλάδας, ώστε να
      ενσωματώνει τις ενημερωμένες προβλέψεις φορτίου και ανανεώσιμης παραγωγής. Στα κείμενα
      του ΑΔΜΗΕ η ΕΔΣ αναφέρεται και ως «εφεδρεία συγκράτησης συχνότητας».
\item \textbf{Αγορά Ενέργειας Εξισορρόπησης}, με τις διαδικασίες χΕΑΣ και αΕΑΣ. Οι οντότητες
      που παρέχουν υπηρεσίες εξισορρόπησης υποβάλλουν προσφορές ενέργειας εξισορρόπησης, τις
      οποίες ο ΑΔΜΗΕ ενεργοποιεί σε πραγματικό χρόνο ανάλογα με την κατάσταση του συστήματος.
\item \textbf{Εκκαθάριση}, στην οποία υπολογίζονται οι αποκλίσεις κάθε συμμετέχοντα από το
      πρόγραμμά του ανά 15λεπτη \emph{Περίοδο Εκκαθάρισης Αποκλίσεων}. Οι αποκλίσεις
      τιμολογούνται με μία Τιμή Αποκλίσεων ανά περίοδο, η οποία προκύπτει ως μέσο κόστος της
      ενέργειας εξισορρόπησης που ενεργοποιήθηκε στην κυρίαρχη κατεύθυνση, ανοδική ή καθοδική.
\end{enumerate}

Η Τιμή Αποκλίσεων δημιουργεί το οικονομικό κίνητρο για ακριβή προγραμματισμό: συμμετέχων με
έλλειμμα σε περίοδο κατά την οποία το σύστημα ενεργοποιεί ανοδική ενέργεια εξισορρόπησης
πληρώνει την ενέργεια που του λείπει στην Τιμή Αποκλίσεων, η οποία αντανακλά το κόστος της
ενέργειας αυτής. Οι κανόνες λειτουργίας περιέχονται στον
Κανονισμό Αγοράς Επόμενης Ημέρας και Ενδοημερήσιας Αγοράς του ΕΧΕ, στον
\href{https://www.admie.gr/agora/rythmistiko-plaisio-agoras/kanonismos-agoras-eksisorropisis}{Κανονισμό
Αγοράς Εξισορρόπησης} και στον
\href{https://www.admie.gr/agora/rythmistiko-plaisio-agoras/kodikas-diaxeirisis-esmie}{Κώδικα
Διαχείρισης του ΕΣΜΗΕ}, οι οποίοι εγκρίνονται από τη Ρυθμιστική Αρχή Αποβλήτων, Ενέργειας και
Υδάτων (ΡΑΑΕΥ).

\textbf{Δεδομένα λειτουργίας και αγοράς σε ηλεκτρονική μορφή.} Οι διαχειριστές δημοσιεύουν
καθημερινά τα αποτελέσματα των αγορών και της λειτουργίας του συστήματος:

\begin{itemize}
\item ΕΧΕ:
      \href{https://www.enexgroup.gr/el/web/guest/markets-publications-el-day-ahead-market}{αποτελέσματα
      της Αγοράς Επόμενης Ημέρας} ανά αγοραία χρονική μονάδα και
      \href{https://www.enexgroup.gr/el/web/guest/day-ahead-market-figures}{διαγράμματα} της ΤΕΑ,
      των όγκων συναλλαγών και των εισαγωγών-εξαγωγών.
\item ΑΔΜΗΕ:
      \href{https://www.admie.gr/agora/statistika-agoras/kyrioi-deiktes-dashboard/paragogi-ana-kaysimo-isozygio-diasyndeseon}{κύριοι
      δείκτες} (παραγωγή ανά καύσιμο, ισοζύγιο διασυνδέσεων, μερίδιο ανανεώσιμων πηγών),
      \href{https://www.admie.gr/agora/genika/agora-eksisorropisis}{περιγραφή της Αγοράς
      Εξισορρόπησης},
      \href{https://www.admie.gr/agora/enimerotika-deltia/ebdomadiaia-deltia-agoras-eksisorropisis}{εβδομαδιαία
      δελτία} της Αγοράς Εξισορρόπησης και
      \href{https://www.admie.gr/agora/statistika-agoras/dedomena}{αρχεία δεδομένων} (φορτίο και
      παραγωγή ανανεώσιμων πηγών από το σύστημα εποπτικού ελέγχου και συλλογής δεδομένων
      (Supervisory Control and Data Acquisition, SCADA), αποτελέσματα των ΔΕΠ, ενέργεια εξισορρόπησης,
      τιμές αποκλίσεων, σφάλμα ελέγχου αποκατάστασης συχνότητας ανά 15 λεπτά), τα οποία
      διατίθενται και μέσω
      \href{https://www.admie.gr/agora/statistika-agoras/file-download-api}{διεπαφής
      προγραμματισμού εφαρμογών} (Application Programming Interface, API) για αυτόματη λήψη.
\end{itemize}

Τα αρχεία του ΑΔΜΗΕ δημοσιεύονται συνήθως την επόμενη ημέρα και αποτελούν κατάλληλη βάση για
ασκήσεις ανάλυσης δεδομένων λειτουργίας του ελληνικού συστήματος.

\section{Καταστάσεις λειτουργίας και κριτήριο N-1}
\label{sec:states}

Η κατάταξη της λειτουργίας σε διακριτές καταστάσεις ασφάλειας, με διαφορετική στρατηγική
ελέγχου σε καθεμία, διατυπώθηκε από τον Dy Liacco στα τέλη της δεκαετίας του 1960. Ο
Κανονισμός (ΕΕ) 2017/1485 για τη λειτουργία του συστήματος μεταφοράς (System Operation
Guideline, SO GL) ορίζει στο άρθρο 18 πέντε καταστάσεις του συστήματος:

\begin{enumerate}
\item \textbf{Κανονική κατάσταση} (normal state). Οι τάσεις, οι ροές ισχύος και η συχνότητα
      βρίσκονται εντός ορίων και το σύστημα παραμένει εντός ορίων μετά από οποιοδήποτε
      απρόβλεπτο συμβάν του καταλόγου απρόβλεπτων συμβάντων, λαμβανομένων υπόψη των
      διαθέσιμων διορθωτικών μέτρων. Ο έλεγχος είναι \emph{προληπτικός} (preventive).
\item \textbf{Κατάσταση συναγερμού} (alert state· συχνά και «κατάσταση επιφυλακής»). Τα
      μεγέθη βρίσκονται εντός ορίων, αλλά τουλάχιστον ένα απρόβλεπτο συμβάν του καταλόγου
      οδηγεί σε παραβίαση ορίων ακόμη και μετά την εφαρμογή διορθωτικών μέτρων, ή η
      εφεδρική δυναμικότητα έχει μειωθεί κατά περισσότερο από 20\% για διάστημα
      μεγαλύτερο των 30 λεπτών.
\item \textbf{Κατάσταση έκτακτης ανάγκης} (emergency state). Έχει παραβιαστεί τουλάχιστον
      ένα όριο επιχειρησιακής ασφάλειας ή εφαρμόζεται μέτρο του σχεδίου άμυνας του
      συστήματος. Ο έλεγχος είναι \emph{διορθωτικός} (corrective).
\item \textbf{Κατάσταση γενικής διακοπής} (blackout state· συχνά και «συσκότιση»). Έχει
      απολεσθεί περισσότερο από το 50\% της ζήτησης ή απουσιάζει πλήρως η τάση για
      τουλάχιστον τρία λεπτά στην περιοχή ελέγχου του ΔΣΜ.
\item \textbf{Κατάσταση αποκατάστασης} (restoration state). Ο ΔΣΜ έχει αρχίσει να
      εφαρμόζει τα μέτρα του σχεδίου αποκατάστασης για την επανατροφοδότηση του
      συστήματος.
\end{enumerate}

Η μετάβαση από την κανονική κατάσταση στην κατάσταση συναγερμού δεν συνοδεύεται από
παραβίαση ορίου ή από μεταβολή των μετρούμενων μεγεθών· προκύπτει μόνο από την ανάλυση
των απρόβλεπτων συμβάντων. Ο προσδιορισμός της κατάστασης του συστήματος προϋποθέτει
επομένως ακριβές μοντέλο της τρέχουσας λειτουργίας.

\subsection{Το κριτήριο N-1}
\label{sec:n1}

Σύμφωνα με το \textbf{κριτήριο (N-1)} του SO GL, τα στοιχεία που παραμένουν σε λειτουργία
εντός της περιοχής ελέγχου του ΔΣΜ μετά την εμφάνιση απρόβλεπτου συμβάντος πρέπει να
εξυπηρετούν τη νέα κατάσταση λειτουργίας χωρίς παραβίαση των ορίων επιχειρησιακής ασφάλειας
και χωρίς αλυσιδωτές αποζεύξεις. Η κατάσταση στην οποία κανένα στοιχείο δεν είναι μη
διαθέσιμο λόγω απρόβλεπτου συμβάντος ονομάζεται \emph{κατάσταση N}, ενώ η κατάσταση μετά την
εμφάνιση συμβάντος του καταλόγου απρόβλεπτων συμβάντων ονομάζεται \emph{κατάσταση (N-1)}. Το
όνομα του κριτηρίου προέρχεται από την απώλεια ενός οποιουδήποτε από τα $N$ στοιχεία του
συστήματος: γραμμής, μετασχηματιστή, μονάδας ηλεκτροπαραγωγής ή φορτίου.

\textbf{Κατάλογος απρόβλεπτων συμβάντων.} Κάθε ΔΣΜ καταρτίζει κατάλογο των εσωτερικών και
εξωτερικών απρόβλεπτων συμβάντων που μπορούν να θέσουν σε κίνδυνο την επιχειρησιακή ασφάλεια
της περιοχής ελέγχου του. Ο SO GL κατατάσσει τα συμβάντα σε τρεις κατηγορίες:

\begin{itemize}
\item \emph{Σύνηθες απρόβλεπτο συμβάν} (ordinary contingency): απρόβλεπτο συμβάν σε έναν
      μόνο κλάδο ή μία μόνο έγχυση, όπως η απώλεια μίας γραμμής, ενός μετασχηματιστή ή μίας
      μονάδας ηλεκτροπαραγωγής.
\item \emph{Σπάνιο απρόβλεπτο συμβάν} (exceptional contingency): ταυτόχρονη εμφάνιση πολλών
      απρόβλεπτων συμβάντων με κοινή αιτία, όπως η απώλεια γραμμής διπλού κυκλώματος σε κοινούς
      πυλώνες ή σφάλμα σε ζυγό υποσταθμού που αποσυνδέει όλες τις αναχωρήσεις του.
\item \emph{Εκτός πεδίου απρόβλεπτο συμβάν} (out-of-range contingency): ταυτόχρονη εμφάνιση
      πολλών απρόβλεπτων συμβάντων χωρίς κοινή αιτία ή απώλεια παραγωγής που υπερβαίνει το
      συμβάν αναφοράς.
\end{itemize}

Ο κατάλογος περιλαμβάνει τα συνήθη απρόβλεπτα συμβάντα και τα σπάνια, στο μέτρο που η
πιθανότητα επέλευσής τους είναι σημαντική, για παράδειγμα όταν τη αυξάνουν οι καιρικές
συνθήκες. Τα εκτός πεδίου απρόβλεπτα συμβάντα δεν καλύπτονται από το κριτήριο· η
αντιμετώπισή τους βασίζεται στα μέτρα του σχεδίου άμυνας του συστήματος, όπως η αυτόματη
αποσύνδεση ζήτησης λόγω χαμηλής συχνότητας.

\textbf{Μαθηματική διατύπωση.} Με $\mathcal{C}$ τον κατάλογο απρόβλεπτων συμβάντων και
$I_l^{(k)}$, $V_i^{(k)}$ το ρεύμα του κλάδου $l$ και το μέτρο της τάσης του ζυγού $i$ στην
κατάσταση που προκύπτει μετά το συμβάν $k$, το κριτήριο απαιτεί

\begin{equation}
I_l^{(k)} \le I_l^{\max}, \qquad V_i^{\min} \le V_i^{(k)} \le V_i^{\max},
\qquad \forall\, l,\ i,\ k \in \mathcal{C}
\label{eq:n1}
\end{equation}

για κάθε κλάδο $l$, κάθε ζυγό $i$ και κάθε συμβάν $k$ του καταλόγου, καθώς και διατήρηση της
ευστάθειας. Στην κατάσταση (N-1) επιτρέπεται η αξιοποίηση
\emph{προσωρινών επιτρεπτών υπερφορτίσεων}, δηλαδή φορτίσεων πάνω από το μόνιμο όριο για
περιορισμένο χρονικό διάστημα, οι οποίες δεν προκαλούν φυσική ζημία στον εξοπλισμό, εφόσον ο
ΔΣΜ έχει εκπονήσει διορθωτικά μέτρα που επαναφέρουν τη φόρτιση εντός του μόνιμου ορίου μέσα
στο διάστημα αυτό.

Το κριτήριο ικανοποιείται με δύο τρόπους:

\begin{itemize}
\item \textbf{Προληπτικά} (preventive security). Το σημείο λειτουργίας της κατάστασης N
      επιλέγεται, με ανακατανομή παραγωγής, αλλαγές τοπολογίας ή ρύθμιση μετασχηματιστών
      μετατόπισης φάσης, ώστε η \eqref{eq:n1} να ισχύει χωρίς καμία ενέργεια μετά το
      συμβάν. Το κόστος των μέτρων καταβάλλεται συνεχώς, ανεξάρτητα από το αν θα εμφανιστεί
      το συμβάν.
\item \textbf{Με διορθωτικά μέτρα μετά το συμβάν} (corrective ή curative security). Στην
      κατάσταση (N-1) γίνεται δεκτή προσωρινή υπέρβαση του μόνιμου ορίου, υπό την
      προϋπόθεση ότι έχουν προετοιμαστεί αυτόματα ή χειροκίνητα μέτρα που την εξαλείφουν
      εγκαίρως. Το κόστος καταβάλλεται μόνο αν εμφανιστεί το συμβάν, αλλά η ασφάλεια
      εξαρτάται από την αξιοπιστία των μέτρων.
\end{itemize}

Ο συνδυασμός τους διατυπώνεται ως βέλτιστη ροή φορτίου με περιορισμούς ασφάλειας
(Security-Constrained Optimal Power Flow, SCOPF), δηλαδή ελαχιστοποίηση του κόστους
λειτουργίας υπό τους περιορισμούς της κατάστασης N και της \eqref{eq:n1} για κάθε συμβάν του
καταλόγου.

\textbf{Ανάλυση απρόβλεπτων συμβάντων.} Η συμμόρφωση ελέγχεται με \emph{ανάλυση απρόβλεπτων
συμβάντων} (contingency analysis· στη βιβλιογραφία και «ανάλυση ενδεχομένων»): για κάθε
συμβάν του καταλόγου εκτελείται υπολογισμός ροής φορτίου (load flow) με το αντίστοιχο
στοιχείο εκτός λειτουργίας και ελέγχονται τα όρια. Για σύστημα μερικών εκατοντάδων κλάδων
απαιτούνται εκατοντάδες υπολογισμοί ροής φορτίου σε κάθε κύκλο ανάλυσης, ο οποίος πρέπει να
ολοκληρώνεται εντός λίγων λεπτών. Για τη μείωση του υπολογιστικού φόρτου εφαρμόζεται
προκαταρκτική επιλογή των κρισιμότερων συμβάντων (contingency screening) με γραμμικοποιημένους
συντελεστές κατανομής. Στο γραμμικό μοντέλο ΣΡ, η ροή του κλάδου $l$ μετά την απώλεια του
κλάδου $k$ δίνεται από τον \emph{συντελεστή κατανομής διακοπής γραμμής} (Line Outage
Distribution Factor, LODF):

\begin{equation}
P_l^{(k)} = P_l^{0} + \mathrm{LODF}_{l,k}\, P_k^{0}, \qquad
\mathrm{LODF}_{l,k} = \frac{\mathrm{PTDF}_{l,k}}{1 - \mathrm{PTDF}_{k,k}}
\label{eq:lodf}
\end{equation}

όπου $P^{0}$ οι ροές στην κατάσταση N και $\mathrm{PTDF}_{l,k}$ ο συντελεστής κατανομής
μεταφοράς ισχύος (Power Transfer Distribution Factor), δηλαδή το ποσοστό μιας μεταφοράς
ισχύος από το ένα άκρο του κλάδου $k$ προς το άλλο που διέρχεται από τον κλάδο $l$. Οι
συντελεστές υπολογίζονται μία φορά για δεδομένη τοπολογία, οπότε η επίδραση όλων των
απωλειών κλάδων εκτιμάται με πράξεις πινάκων· μόνο τα συμβάντα που οδηγούν σε φόρτιση κοντά
στο όριο εξετάζονται στη συνέχεια με πλήρη υπολογισμό ροής φορτίου εναλλασσόμενου ρεύματος.

\emph{Παράδειγμα.} Δύο όμοιες παράλληλες γραμμές, με όριο 1\,000 MW η καθεμία, αποτελούν τη
μοναδική σύνδεση μεταξύ δύο περιοχών. Στην κατάσταση N η μεταφορά μοιράζεται εξίσου, οπότε
$\mathrm{PTDF}_{k,k} = \mathrm{PTDF}_{l,k} = 0{,}5$ και από την \eqref{eq:lodf}
$\mathrm{LODF}_{l,k} = 1$: με την απώλεια της μίας γραμμής, η άλλη αναλαμβάνει ολόκληρη τη
μεταφορά. Το κριτήριο N-1 περιορίζει επομένως τη συνολική μεταφορά σε 1\,000 MW, δηλαδή σε
φόρτιση κάθε γραμμής έως 50\% στην κατάσταση N, αν και η θερμική ικανότητα της σύνδεσης είναι
2\,000 MW. Αν η γραμμή επιτρέπει προσωρινή υπερφόρτιση 20\% και έχουν προετοιμαστεί
διορθωτικά μέτρα, όπως ανακατανομή παραγωγής που εφαρμόζεται εντός του επιτρεπόμενου χρόνου,
η μεταφορά μπορεί να αυξηθεί έως 1\,200 MW.

Το κριτήριο N-1 είναι \emph{ντετερμινιστικό}: δεν λαμβάνει υπόψη την πιθανότητα εμφάνισης
κάθε συμβάντος και αντιμετωπίζει ισοδύναμα τα σπάνια και τα συχνά συμβάντα. Είναι απλό στην
εφαρμογή και στον έλεγχο, αλλά σε ορισμένες συνθήκες οδηγεί σε συντηρητική λειτουργία και σε
άλλες δεν καλύπτει πολλαπλά ή συσχετισμένα συμβάντα· για τον λόγο αυτόν αναπτύσσονται
πιθανοτικά κριτήρια αξιοπιστίας, τα οποία σταθμίζουν κάθε συμβάν με την πιθανότητα και τις
συνέπειές του. Η ανάλυση απρόβλεπτων συμβάντων χρησιμοποιεί ως αρχική συνθήκη την τρέχουσα
κατάσταση του συστήματος, η οποία παρέχεται από την εκτίμηση κατάστασης· σφάλμα της
εκτίμησης μεταφέρεται απευθείας στην αξιολόγηση της ασφάλειας.

\section{Εποπτεία και έλεγχος: SCADA και EMS}

\subsection{Αλυσίδα μέτρησης}

Η υποδομή τηλεμέτρησης και τηλεχειρισμού του ΣΜΗΕ βασίζεται στο σύστημα εποπτικού ελέγχου
και συλλογής δεδομένων (Supervisory Control and Data Acquisition, \textbf{SCADA}). Η αλυσίδα
από το φυσικό μέγεθος έως το κέντρο ελέγχου περιλαμβάνει τα ακόλουθα στοιχεία:

\begin{enumerate}
\item \textbf{Μετασχηματιστές μέτρησης.} Οι μετασχηματιστές τάσης (Voltage Transformers, VT) και
      έντασης (Current Transformers, CT)
      υποβιβάζουν τα μεγέθη του δικτύου σε τιμές κατάλληλες για τις διατάξεις μέτρησης
      και προστασίας, τυπικά 100 V και 1 ή 5 A στο δευτερεύον τύλιγμα. Εισάγουν το πρώτο
      σφάλμα της αλυσίδας, το οποίο καθορίζεται από την κλάση ακρίβειας (τυπικά 0,2 έως 1,
      δηλαδή σφάλμα λόγου 0,2 έως 1\%).
\item \textbf{Μετατροπείς μέτρησης (transducers) και ευφυείς ηλεκτρονικές συσκευές (Intelligent
      Electronic Devices, IED).}
      Δειγματοληπτούν τις κυματομορφές και υπολογίζουν τα παράγωγα μεγέθη: ενεργό και
      άεργο ισχύ, ενεργό τιμή (root mean square, rms) τάσης και έντασης.
\item \textbf{Απομακρυσμένες τερματικές μονάδες (Remote Terminal Units, RTU).} Συγκεντρώνουν τα δεδομένα του
      υποσταθμού, τα μεταδίδουν στο Κέντρο Ελέγχου Ενέργειας και εκτελούν τις εντολές
      τηλεχειρισμού.
\item \textbf{Τηλεπικοινωνιακό δίκτυο.} Μεταφέρει τα δεδομένα με τυποποιημένα πρωτόκολλα,
      κυρίως IEC 60870-5-101/104 και DNP3 (Distributed Network Protocol 3)· εντός του υποσταθμού χρησιμοποιείται σήμερα
      κατά κανόνα το IEC 61850.
\end{enumerate}

Τα δεδομένα διακρίνονται σε δύο κατηγορίες:

\begin{itemize}
\item \textbf{Αναλογικές μετρήσεις} (τηλεμετρήσεις): ροές ενεργού και αέργου ισχύος στα
      άκρα των γραμμών και των μετασχηματιστών, εγχύσεις ισχύος στους ζυγούς, μέτρα τάσεων
      ζυγών, εντάσεις ρεύματος κλάδων.
\item \textbf{Ψηφιακές ενδείξεις} (τηλεσημάνσεις): θέσεις διακοπτών ισχύος και αποζευκτών,
      θέσεις των μεταγωγέων λήψεων υπό φορτίο των μετασχηματιστών, σήματα συναγερμού και
      λειτουργίας των ηλεκτρονόμων προστασίας. Οι ενδείξεις θέσης καθορίζουν την
      \emph{τοπολογία} του δικτύου, δηλαδή τα στοιχεία που βρίσκονται σε λειτουργία και τον
      τρόπο σύνδεσής τους.
\end{itemize}

Η περίοδος σάρωσης ενός συμβατικού συστήματος SCADA είναι τυπικά 2 έως 10 s.

\subsection{Περιορισμοί των μετρήσεων SCADA}

Οι μετρήσεις SCADA παρουσιάζουν δύο περιορισμούς, οι οποίοι καθορίζουν τις απαιτήσεις της
εκτίμησης κατάστασης:

\begin{itemize}
\item \textbf{Απουσία χρονικού συγχρονισμού.} Οι μετρήσεις λαμβάνονται σε διαφορετικές
      χρονικές στιγμές εντός του κύκλου σάρωσης και δεν αντιστοιχούν σε κοινό χρονικό
      στιγμιότυπο. Το σύνολό τους αποτελεί \emph{οιονεί στατική} (quasi-static)
      αναπαράσταση του συστήματος, αποδεκτή σε μόνιμη κατάσταση λειτουργίας αλλά
      ανακριβή κατά τη διάρκεια μεταβατικών φαινομένων.
\item \textbf{Θόρυβος, εσφαλμένα δεδομένα και ελλιπής κάλυψη.} Οι μετρήσεις περιέχουν
      τυχαία σφάλματα και ενίοτε χονδροειδή σφάλματα (gross errors), ενώ ορισμένα μεγέθη
      δεν μετρώνται. Οι γωνίες των τάσεων των ζυγών, οι οποίες αποτελούν το ήμισυ των
      μεταβλητών κατάστασης, δεν μετρώνται από το συμβατικό SCADA.
\end{itemize}

Χονδροειδή σφάλματα προκαλούνται, μεταξύ άλλων, από εσφαλμένα καταχωρημένο λόγο
μετασχηματισμού μετασχηματιστή έντασης, από ανεστραμμένη πολικότητα μέτρησης ισχύος ή από
ένδειξη θέσης διακόπτη που δεν ανανεώνεται λόγω βλάβης. Τα δεδομένα αυτά δεν διακρίνονται
από τα ορθά με έλεγχο κάθε μέτρησης χωριστά· ανιχνεύονται μόνο με σύγκριση των μετρήσεων
μεταξύ τους, με αξιοποίηση του πλεονασμού των μετρήσεων (redundancy).

\subsection{Λειτουργίες του EMS}

Το σύστημα διαχείρισης ενέργειας (Energy Management System, \textbf{EMS}) είναι το σύνολο των
εφαρμογών λογισμικού του κέντρου ελέγχου που επεξεργάζονται τα δεδομένα του SCADA. Οι
βασικές λειτουργίες εκτελούνται με την ακόλουθη σειρά, καθώς η έξοδος καθεμίας αποτελεί
είσοδο της επόμενης:

\begin{enumerate}
\item \textbf{Επεξεργαστής τοπολογίας} (topology processor). Προσδιορίζει, από τις θέσεις
      των διακοπτών και των αποζευκτών, την τρέχουσα τοπολογία και σχηματίζει το μοντέλο
      ζυγών-κλάδων (bus-branch model) του δικτύου.
\item \textbf{Εκτίμηση κατάστασης}. Προσδιορίζει τις τάσεις (μέτρα και γωνίες) όλων των
      ζυγών από τις διαθέσιμες μετρήσεις.
\item \textbf{Ανάλυση απρόβλεπτων συμβάντων}. Αξιολογεί την επιχειρησιακή ασφάλεια της
      εκτιμώμενης κατάστασης έναντι του καταλόγου απρόβλεπτων συμβάντων.
\item \textbf{Βέλτιστη ροή φορτίου} (Optimal Power Flow, OPF). Προσδιορίζει, εφόσον
      απαιτείται, τα αναγκαία διορθωτικά μέτρα (remedial actions) με το ελάχιστο κόστος.
\item \textbf{Αυτόματη Ρύθμιση Παραγωγής} (ΑΡΠ· Automatic Generation Control, AGC).
      Υλοποιεί συνεχώς τη ρύθμιση φορτίου-συχνότητας και την τήρηση των προγραμματισμένων
      ανταλλαγών ισχύος με τις γειτονικές περιοχές (Κεφάλαιο 3).
\end{enumerate}

Σφάλμα του επεξεργαστή τοπολογίας ή της εκτίμησης κατάστασης μεταφέρεται σε όλες τις
επόμενες λειτουργίες χωρίς να γίνεται αντιληπτό.

Οι χειριστές του κέντρου ελέγχου εποπτεύουν το σύστημα μέσω σταθμών εργασίας και τοίχων
προβολής (Σχήμα~\ref{fig:control-room}). Στους σταθμούς εργασίας εμφανίζονται μονογραμμικά
διαγράμματα των υποσταθμών με τις θέσεις των διακοπτών, τις ροές ισχύος και τις τάσεις,
καμπύλες φορτίου και πρόβλεψης, καθώς και οι κατάλογοι συναγερμών του SCADA· από τους ίδιους
σταθμούς εκδίδονται οι εντολές τηλεχειρισμού. Οι τοίχοι προβολής δίνουν συνοπτική εικόνα
ολόκληρου του δικτύου, ώστε όλοι οι χειριστές να έχουν κοινή αντίληψη της κατάστασης του
συστήματος.

\begin{figure}[!htbp]
\centering
\includegraphics[width=\linewidth]{figures/terna-kentro-elegchou.jpg}
\caption{Αίθουσα ελέγχου κατανομής της Terna, Διαχειριστή Συστήματος Μεταφοράς της Ιταλίας.
Στους σταθμούς εργασίας εμφανίζονται μονογραμμικά διαγράμματα και καμπύλες φορτίου, στον τοίχο
προβολής συνοπτικές απεικονίσεις του δικτύου και αριστερά μονογραμμικό διάγραμμα του δικτύου
υπερυψηλής τάσης. Φωτογραφία: Terna S.p.A.,
\href{https://commons.wikimedia.org/wiki/File:Dispacciamento_-_Sala_controllo.JPG}{«Dispacciamento
- Sala controllo»}, άδεια
\href{https://creativecommons.org/licenses/by-sa/3.0/}{CC BY-SA 3.0}, μέσω Wikimedia Commons.}
\label{fig:control-room}
\end{figure}

Το Σχήμα~\ref{fig:scada-ems} απεικονίζει τον κλειστό βρόχο εποπτείας και ελέγχου: οι
μετρήσεις μεταδίδονται από το δίκτυο προς το κέντρο ελέγχου και οι εντολές ελέγχου
επιστρέφουν προς το δίκτυο.

\begin{figure}[!htbp]
\centering
\includegraphics[width=\linewidth]{figures/scada-ems.svg}
\caption{Ο βρόχος εποπτείας και ελέγχου. Η επάνω σειρά είναι η αλυσίδα μέτρησης, από το
φυσικό δίκτυο έως το SCADA· η κάτω σειρά είναι η ακολουθία λειτουργιών του EMS, η οποία
καταλήγει σε εντολές προς το δίκτυο.}
\label{fig:scada-ems}
\end{figure}

\section{Συγχρονισμένες μετρήσεις φασιθετών και συστήματα WAMS}

\subsection{Γωνία τάσης και ροή ενεργού ισχύος}

Η ενεργός ισχύς που μεταφέρεται μεταξύ δύο ζυγών μέσω γραμμής χωρίς απώλειες με επαγωγική
αντίδραση σειράς $X$ δίνεται από τη σχέση

\begin{equation}
P_{12} = \frac{V_1 V_2}{X}\,\sin(\delta_1 - \delta_2)
\label{eq:powerflow}
\end{equation}

Για τάσεις κοντά στην ονομαστική τιμή, η ροή ενεργού ισχύος καθορίζεται κυρίως από τη
διαφορά γωνιών $\delta_1 - \delta_2$, η οποία αποτελεί άμεσο μέτρο της φόρτισης της γραμμής
και του περιθωρίου στατικής ευστάθειας· η μέγιστη μεταφερόμενη ισχύς αντιστοιχεί σε
$\delta_1 - \delta_2 = 90^\circ$. Το συμβατικό SCADA δεν μετρά γωνίες τάσης.

\subsection{Αρχή λειτουργίας της PMU}

Οι \textbf{μονάδες μέτρησης φασιθετών} (Phasor Measurement Units, PMU) υπολογίζουν τους
φασιθέτες — μέτρο και γωνία — της τάσης και της έντασης του ρεύματος, ανά φάση και για τη
θετική ακολουθία, με χρονοσήμανση ως προς την κοινή χρονική αναφορά της συντονισμένης
παγκόσμιας ώρας (Coordinated Universal Time, UTC), η οποία λαμβάνεται τυπικά από δέκτη του
δορυφορικού συστήματος εντοπισμού θέσης (Global Positioning System, GPS).

Η γωνία ενός φασιθέτη ορίζεται ως προς σήμα αναφοράς. Επειδή όλες οι PMU χρησιμοποιούν την
ίδια χρονική βάση, οι μετρούμενες γωνίες είναι άμεσα συγκρίσιμες σε ολόκληρο το σύστημα,
ανεξάρτητα από τη γεωγραφική απόσταση των σημείων μέτρησης· οι μετρήσεις αυτές ονομάζονται
\emph{συγχρονισμένοι φασιθέτες} (synchrophasors). Η απαιτούμενη ακρίβεια χρονισμού είναι της
τάξης του μικροδευτερολέπτου: σε σύστημα 50 Hz, σφάλμα χρονισμού 1 μs αντιστοιχεί σε σφάλμα
γωνίας $360^\circ \times 50 \times 10^{-6}$, δηλαδή 0,018°.

\begin{figure}[!htbp]
\centering
\includegraphics[width=\linewidth]{figures/pmu-wams.svg}
\caption{Αριστερά, το μετρούμενο μέγεθος: η γωνία $\delta$ της κυματομορφής τάσης ορίζεται
ως προς τον παλμό ανά δευτερόλεπτο (pulse per second, PPS) της κοινής χρονικής αναφοράς UTC και μαζί με το
μέτρο σχηματίζει τον φασιθέτη. Δεξιά, η ιεραρχία συλλογής: οι PMU συγχρονίζονται μέσω GPS, τα
πλαίσια δεδομένων ευθυγραμμίζονται χρονικά στον συγκεντρωτή δεδομένων φασιθετών (Phasor Data
Concentrator, PDC) και τροφοδοτούν το σύστημα εποπτείας ευρείας περιοχής (Wide Area
Monitoring System, WAMS).}
\label{fig:pmu-wams}
\end{figure}

\subsection{Υπολογισμός του φασιθέτη}
\label{sec:phasor-estimation}

\textbf{Ορισμός του συγχρονισμένου φασιθέτη.} Ημιτονοειδές σήμα ονομαστικής συχνότητας

\begin{equation}
x(t) = \sqrt{2}\,X \cos(\omega_0 t + \varphi), \qquad \omega_0 = 2\pi f_0
\label{eq:sinusoid}
\end{equation}

αναπαρίσταται από τον φασιθέτη $\tilde{X} = X e^{j\varphi} = X\cos\varphi + jX\sin\varphi$,
όπου $X$ η ενεργός τιμή. Στον \emph{συγχρονισμένο} φασιθέτη η γωνία $\varphi$ μετράται ως
προς συνημίτονο ονομαστικής συχνότητας συγχρονισμένο με την UTC: ισχύει $\varphi = 0$ όταν η
κορυφή του σήματος συμπίπτει με την αλλαγή του δευτερολέπτου UTC (παλμός PPS, Σχήμα
\ref{fig:pmu-wams}) και $\varphi = -90^\circ$ όταν συμπίπτει η θετική διέλευση από το μηδέν.

Όταν η συχνότητα αποκλίνει από την ονομαστική, $f = f_0 + \Delta f$, το σήμα γράφεται
$x(t) = \sqrt{2}\,X \cos\left(\omega_0 t + 2\pi\Delta f\,t + \varphi\right)$ και ο φασιθέτης
στρέφεται με ταχύτητα ίση με τη διαφορά συχνότητας:

\begin{equation}
\tilde{X}(t) = X e^{j\phi(t)}, \qquad \phi(t) = 2\pi\Delta f\,t + \varphi
\label{eq:rotphasor}
\end{equation}

Η συχνότητα και ο ρυθμός μεταβολής της συχνότητας, που αναφέρει επίσης η PMU, προκύπτουν
από τις παραγώγους της γωνίας:

\begin{equation}
f(t) = f_0 + \frac{1}{2\pi}\,\frac{\mathrm{d}\phi}{\mathrm{d}t}, \qquad
\mathrm{RoCoF}(t) = \frac{\mathrm{d}f}{\mathrm{d}t}
= \frac{1}{2\pi}\,\frac{\mathrm{d}^2\phi}{\mathrm{d}t^2}
\label{eq:pmufreq}
\end{equation}

\textbf{Εκτίμηση με διακριτό μετασχηματισμό Fourier.} Η PMU δειγματοληπτεί το σήμα με
συχνότητα $f_s = N f_0$, δηλαδή με $N$ δείγματα ανά ονομαστικό κύκλο, και εκτιμά τον
φασιθέτη από τη θεμελιώδη συνιστώσα του διακριτού μετασχηματισμού Fourier (Discrete Fourier
Transform, DFT) σε παράθυρο ενός κύκλου:

\begin{equation}
\hat{\tilde{X}} = \frac{\sqrt{2}}{N} \sum_{k=0}^{N-1} x_k\, e^{-j 2\pi k / N},
\qquad x_k = x(k/f_s)
\label{eq:dft}
\end{equation}

Για σήμα ονομαστικής συχνότητας η εκτίμηση είναι ακριβής. Με τη βοήθεια της ταυτότητας
Euler τα δείγματα γράφονται
$x_k = \frac{X}{\sqrt{2}}\left[e^{j(2\pi k/N + \varphi)} + e^{-j(2\pi k/N + \varphi)}\right]$.
Μετά τον πολλαπλασιασμό με $e^{-j2\pi k/N}$, ο πρώτος όρος γίνεται σταθερός και δίνει
άθροισμα $N X e^{j\varphi}/\sqrt{2}$, ενώ ο δεύτερος μηδενίζεται, επειδή
$\sum_{k=0}^{N-1} e^{-j4\pi k/N} = 0$ για $N > 2$· επομένως $\hat{\tilde{X}} = X e^{j\varphi}$.
Η γωνία αναφέρεται στη χρονική στιγμή του πρώτου δείγματος, η οποία είναι γνωστή από τη
χρονοσήμανση· στην πράξη η μέτρηση αποδίδεται στο κέντρο του παραθύρου και η γωνία
διορθώνεται αντίστοιχα.

Σε ολισθαίνον παράθυρο η εκτίμηση ανανεώνεται με κάθε νέο δείγμα. Αν $\hat{\tilde{X}}^{(r)}$
είναι η εκτίμηση από τα δείγματα $x_r, \ldots, x_{r+N-1}$, με την ίδια απόλυτη αρίθμηση
$k$ στον εκθέτη της \eqref{eq:dft}, ισχύει η αναδρομική σχέση

\begin{equation}
\hat{\tilde{X}}^{(r+1)} = \hat{\tilde{X}}^{(r)}
+ \frac{\sqrt{2}}{N}\left(x_{r+N} - x_r\right) e^{-j2\pi r/N}
\label{eq:recdft}
\end{equation}

η οποία απαιτεί μία μόνο πράξη ανά δείγμα. Για σήμα ονομαστικής συχνότητας η εκτίμηση
παραμένει σταθερή, ενώ για $f \neq f_0$ στρέφεται σύμφωνα με την \eqref{eq:rotphasor}.

\textbf{Σφάλματα εκτός ονομαστικής συχνότητας.} Όταν $f \neq f_0$, το παράθυρο δεν περιέχει
ακέραιο αριθμό περιόδων και ο δεύτερος όρος δεν μηδενίζεται πλήρως (διαρροή φάσματος,
spectral leakage). Η εκτίμηση περιέχει τότε μικρό σφάλμα μέτρου και γωνίας και ταλάντωση
συχνότητας περίπου $2f_0$. Για τον λόγο αυτόν οι PMU εφαρμόζουν ψηφιακά φίλτρα: η κλάση P
χρησιμοποιεί σύντομο παράθυρο, στο πρότυπο αναφοράς τριγωνικό μήκους περίπου δύο κύκλων, για
μικρή καθυστέρηση, ενώ η κλάση M μακρύτερα φίλτρα για καλύτερη απόρριψη παρεμβολών.
Επιπλέον, η δειγματοληψία μπορεί να προσαρμόζεται στη μετρούμενη συχνότητα (frequency
tracking).

\textbf{Θετική ακολουθία.} Από τους φασιθέτες των τριών φάσεων υπολογίζεται ο φασιθέτης θετικής
ακολουθίας

\begin{equation}
\tilde{V}_1 = \frac{1}{3}\left(\tilde{V}_a + a\,\tilde{V}_b + a^2\,\tilde{V}_c\right),
\qquad a = e^{j120^\circ}
\label{eq:posseq}
\end{equation}

ο οποίος χρησιμοποιείται στην εκτίμηση κατάστασης και στις εφαρμογές WAMS.


\subsection{Προδιαγραφές}

Τα τεχνικά χαρακτηριστικά των PMU καθορίζονται από το πρότυπο IEC/IEEE 60255-118-1:2018, το
οποίο αντικατέστησε τη σειρά IEEE C37.118:

\begin{itemize}
\item \textbf{Ρυθμός αναφοράς} (reporting rate): τυπικά 10, 25 ή 50 πλαίσια δεδομένων ανά
      δευτερόλεπτο σε σύστημα 50 Hz· επιτρέπονται και υψηλότεροι ρυθμοί.
\item \textbf{Ακρίβεια}: το ολικό διανυσματικό σφάλμα (Total Vector Error, TVE) δεν
      επιτρέπεται να υπερβαίνει το 1\% σε μόνιμη κατάσταση. Το TVE ορίζεται ως το μέτρο της
      διαφοράς μετρούμενου και πραγματικού φασιθέτη, κανονικοποιημένο ως προς το μέτρο του
      πραγματικού, και συνδυάζει σε ένα μέγεθος τα σφάλματα μέτρου, γωνίας και χρονισμού.
      Ορίζονται επιπλέον το σφάλμα συχνότητας (Frequency Error, FE) και το σφάλμα ρυθμού
      μεταβολής συχνότητας (Rate of change of Frequency Error, RFE).
\end{itemize}

Με $X_r$, $X_i$ το πραγματικό και το φανταστικό μέρος του πραγματικού φασιθέτη και
$\hat{X}_r$, $\hat{X}_i$ τα αντίστοιχα της μέτρησης, το ολικό διανυσματικό σφάλμα είναι

\begin{equation}
\mathrm{TVE} = \sqrt{\frac{(\hat{X}_r - X_r)^2 + (\hat{X}_i - X_i)^2}{X_r^2 + X_i^2}}
\label{eq:tve}
\end{equation}

Για σφάλμα μόνο γωνίας $\Delta\varphi$ ισχύει
$\mathrm{TVE} = 2\left|\sin(\Delta\varphi/2)\right| \approx |\Delta\varphi|$. Το όριο 1\%
αντιστοιχεί επομένως σε σφάλμα γωνίας 0,01 rad, δηλαδή 0,573°, και σε σύστημα 50 Hz σε σφάλμα
χρονισμού $0{,}01 / (2\pi \times 50)$ s, δηλαδή περίπου 31,8 μs.

\begin{itemize}
\item \textbf{Κλάσεις}: η κλάση P (protection) προορίζεται για εφαρμογές που απαιτούν ταχεία
      απόκριση και μικρή καθυστέρηση φίλτρου, ενώ η κλάση M (measurement) για εφαρμογές
      που απαιτούν υψηλή ακρίβεια παρουσία αρμονικών και παρεμβολών. Η διάκριση απορρέει
      από τον συμβιβασμό μεταξύ ταχύτητας απόκρισης και απόρριψης παρεμβολών του
      ψηφιακού φίλτρου.
\end{itemize}

\begin{table}[!htbp]
\centering
\begin{tabular}{lll}
\hline
 & SCADA & PMU / WAMS \\
\hline
Ανανέωση         & ανά 2--10 s        & 10--50 πλαίσια ανά s \\
Συγχρονισμός     & όχι                & ναι (UTC μέσω GPS) \\
Γωνία τάσης      & δεν μετράται       & μετράται άμεσα \\
Σχέση με την κατάσταση & μη γραμμική & γραμμική \\
Κάλυψη           & εκτεταμένη         & επιλεκτική, λόγω κόστους \\
Τυπική χρήση     & εκτίμηση κατάστασης, τηλεχειρισμοί & εποπτεία δυναμικών φαινομένων \\
\hline
\end{tabular}
\caption{Σύγκριση μετρήσεων SCADA και PMU. Οι δύο τεχνολογίες χρησιμοποιούνται
συμπληρωματικά.}
\label{tab:scada-pmu}
\end{table}

\subsection{Συστήματα WAMS και εφαρμογές}

Οι μετρήσεις των PMU συγκεντρώνονται ιεραρχικά σε συγκεντρωτές δεδομένων φασιθετών (Phasor
Data Concentrators, PDC), οι οποίοι ευθυγραμμίζουν χρονικά τα πλαίσια δεδομένων και τα
διαθέτουν στο \textbf{σύστημα εποπτείας ευρείας περιοχής} (Wide Area Monitoring System,
WAMS). Οι κύριες εφαρμογές είναι:

\begin{itemize}
\item η παρακολούθηση ηλεκτρομηχανικών ταλαντώσεων μεταξύ περιοχών (inter-area
      oscillations), συχνότητας 0,1 έως 1 Hz, οι οποίες δεν αποτυπώνονται με περίοδο
      σάρωσης δευτερολέπτων·
\item η μέτρηση του ρυθμού μεταβολής της συχνότητας και η εκτίμηση της διαθέσιμης
      αδράνειας·
\item η ανίχνευση νησιδοποίησης, δηλαδή της μετάβασης τμήματος του δικτύου σε απομονωμένη
      λειτουργία, και η επιτήρηση των γωνιακών διαφορών μεταξύ περιοχών·
\item η επικύρωση δυναμικών μοντέλων μετά από διαταραχές, με σύγκριση της καταγεγραμμένης
      και της προσομοιωμένης απόκρισης·
\item η υποστήριξη της εκτίμησης κατάστασης και, σε συνθήκες πλήρους κάλυψης, η
      γραμμική εκτίμηση κατάστασης μόνο από μετρήσεις PMU.
\end{itemize}

\textbf{Συχνότητα σε πραγματικό χρόνο.} Η συχνότητα της συγχρονισμένης περιοχής της
Ηπειρωτικής Ευρώπης δημοσιεύεται σε πραγματικό χρόνο από ανεξάρτητα δίκτυα μετρητών, όπως το
\href{https://gridradar.net/en/mains-frequency}{gridradar} και τα
\href{https://www.energy-charts.info/charts/frequency/chart.htm?l=en}{Energy-Charts} του
Fraunhofer ISE. Στα διαγράμματα αυτά διακρίνονται οι συνεχείς μικρές διακυμάνσεις γύρω από τα
50 Hz, οι απότομες μεταβολές στις αλλαγές των περιόδων προγραμματισμού, όταν τα προγράμματα
παραγωγής μεταβάλλονται βηματικά, και, σπανιότερα, οι διαταραχές από απώλεια μονάδων
ηλεκτροπαραγωγής.

Το Σχήμα~\ref{fig:freq-day} δείχνει τις μετρήσεις ανά 1 s μιας ημέρας, από τα ανοιχτά
δεδομένα του Energy-Charts. Η τυπική απόκλιση της συχνότητας ήταν 24 mHz και η συχνότητα
βρέθηκε εκτός της τυπικής περιοχής $\pm 50$ mHz για 54 λεπτά συνολικά· το όριο του SO GL για
την Ηπειρωτική Ευρώπη, 15\,000 λεπτά ανά έτος, αντιστοιχεί κατά μέσο όρο σε 41 λεπτά ημερησίως.
Η μεγαλύτερη απόκλιση της ημέρας, 110 mHz, εμφανίστηκε αμέσως μετά την αλλαγή του τετάρτου
στις 10:15. Η συμπεριφορά είναι συστηματική: η μέση απόλυτη απόκλιση κορυφώνεται 30 έως 60 s
μετά την αρχή κάθε τετάρτου, στα 29 mHz περίπου, έναντι 15 έως 16 mHz στο μέσο του τετάρτου.
Οι αποκλίσεις αυτές οφείλονται στη βηματική μεταβολή των προγραμμάτων παραγωγής στα όρια των
αγοραίων χρονικών μονάδων (υποενότητα «\nameref{sec:mtu}»), ενώ η ζήτηση μεταβάλλεται συνεχώς.
Το Σχήμα~\ref{fig:freq-live} ενσωματώνει τη μέτρηση της συχνότητας σε πραγματικό χρόνο, στην
οποία οι μεταβολές αυτές παρατηρούνται άμεσα γύρω από κάθε αλλαγή τετάρτου.

\input{embeds/frequency-live.md}

\begin{figure}[!htbp]
\centering
\includegraphics[width=\linewidth]{figures/frequency-day.svg}
\caption{Συχνότητα της Ηπειρωτικής Ευρώπης στις 15 Σεπτεμβρίου 2026, από μετρήσεις ανά 1 s στο
Freiburg της Γερμανίας. (α) Ελάχιστη και μέγιστη τιμή ανά λεπτό· η σκιασμένη ζώνη είναι η
τυπική περιοχή συχνότητας $\pm 50$ mHz. (β) Η μεγαλύτερη απόκλιση της ημέρας, αμέσως μετά την
αλλαγή του τετάρτου στις 10:15. (γ) Μέση απόλυτη απόκλιση ανά 30 s μέσα στο τέταρτο, για όλα
τα τέταρτα της ημέρας. Δεδομένα: \href{https://www.energy-charts.info}{energy-charts.info}
(Fraunhofer ISE), άδεια CC BY 4.0.}
\label{fig:freq-day}
\end{figure}

\section{Εκτίμηση κατάστασης}

\subsection{Διατύπωση του προβλήματος}

Η \textbf{εκτίμηση κατάστασης} (state estimation) προσδιορίζει την πλέον πιθανή κατάσταση
λειτουργίας του δικτύου από ένα πλεονάζον σύνολο μετρήσεων με θόρυβο, με βάση το μοντέλο
δικτύου που σχηματίζει ο επεξεργαστής τοπολογίας. Εισήχθη στα συστήματα ηλεκτρικής ενέργειας
από τους Schweppe και Wildes το 1970 και αποτελεί βασική λειτουργία κάθε EMS.

Ο υπολογισμός ροής φορτίου δεν είναι κατάλληλος για τον σκοπό αυτόν, για δύο λόγους.
Πρώτον, απαιτεί ακριβώς τόσα δεδομένα εισόδου όσοι είναι οι άγνωστοι, ενώ οι διαθέσιμες
μετρήσεις είναι περισσότερες και περιέχουν σφάλματα. Δεύτερον, δεν διαθέτει μηχανισμό
ανίχνευσης εσφαλμένων δεδομένων. Η εκτίμηση κατάστασης αξιοποιεί τον πλεονασμό των
μετρήσεων τόσο για τον περιορισμό της επίδρασης του θορύβου όσο και για την ανίχνευση και
την απόρριψη εσφαλμένων μετρήσεων.

Στη γενική της μορφή, η εκτίμηση κατάστασης είναι πρόβλημα μη γραμμικής βελτιστοποίησης:
οι μετρήσεις συνδέονται με την κατάσταση μέσω των μη γραμμικών εξισώσεων του δικτύου
εναλλασσόμενου ρεύματος. Το Σχήμα~\ref{fig:se-flow} δίνει το διάγραμμα ροής του αλγορίθμου,
ο οποίος περιλαμβάνει δύο εμφωλευμένους βρόχους: τον εσωτερικό, για την επαναληπτική επίλυση
του προβλήματος ελαχιστοποίησης, και τον εξωτερικό, για την απόρριψη εσφαλμένων μετρήσεων.
Σε ένα EMS ο αλγόριθμος εκτελείται περιοδικά, τυπικά ανά ένα έως λίγα λεπτά, καθώς και μετά
από κάθε μεταβολή της τοπολογίας.

\begin{figure}[!htbp]
\centering
\includegraphics[width=0.8\linewidth]{figures/se-flowchart.svg}
\caption{Ροή της εκτίμησης κατάστασης σε ένα EMS, με αριθμημένα βήματα. Ο εσωτερικός βρόχος (αριστερά) επαναλαμβάνει
τη μέθοδο Gauss--Newton έως τη σύγκλιση· ο εξωτερικός (δεξιά) απορρίπτει τη μέτρηση με το
μέγιστο κανονικοποιημένο υπόλοιπο και επανεκτιμά.}
\label{fig:se-flow}
\end{figure}

Τα βήματα του αλγορίθμου, με την αρίθμηση του Σχήματος~\ref{fig:se-flow}, είναι τα ακόλουθα.
Οι επόμενες υποενότητες αναπτύσσουν τη θεωρία κάθε βήματος.

\begin{enumerate}
\item \textbf{Δεδομένα εισόδου.} Οι αναλογικές μετρήσεις SCADA και PMU, οι ψηφιακές ενδείξεις
      θέσης διακοπτών και αποζευκτών και οι παράμετροι των στοιχείων του δικτύου (αντιστάσεις,
      αγωγιμότητες, λόγοι μετασχηματισμού) από τη βάση δεδομένων του EMS. Κάθε μέτρηση
      συνοδεύεται από την τυπική απόκλισή της.
\item \textbf{Επεξεργαστής τοπολογίας.} Στο μοντέλο κόμβων-διακοπτών (node-breaker model)
      κάθε υποσταθμός περιγράφεται με τους διακόπτες και τους αποζεύκτες του. Ο επεξεργαστής
      συγχωνεύει τους κόμβους που συνδέονται με κλειστούς διακόπτες σε ζυγούς, αφαιρεί τα
      στοιχεία εκτός λειτουργίας και σχηματίζει το μοντέλο ζυγών-κλάδων και τον
      $\mathbf{Y}_{\mathrm{bus}}$. Εσφαλμένη ένδειξη θέσης διακόπτη δημιουργεί \emph{σφάλμα
      τοπολογίας}, το οποίο εμφανίζεται ως ομάδα ασυμβίβαστων μετρήσεων γύρω από το
      αντίστοιχο στοιχείο.
\item \textbf{Έλεγχος παρατηρησιμότητας} (υποενότητα «\nameref{sec:observability}»). Αν το σύστημα δεν είναι
      παρατηρήσιμο, προσδιορίζονται οι παρατηρήσιμες νησίδες και προστίθενται
      ψευδομετρήσεις, ώστε να αποκατασταθεί η παρατηρησιμότητα με ελάχιστη επίδραση στην
      εκτίμηση.
\item \textbf{Αρχικοποίηση.} Ως αρχική εκτίμηση $\mathbf{x}^{0}$ χρησιμοποιείται το επίπεδο
      προφίλ ή, σε συνεχή λειτουργία, η λύση της προηγούμενης εκτέλεσης (θερμή εκκίνηση,
      warm start), η οποία μειώνει τον αριθμό των επαναλήψεων.
\item \textbf{Επανάληψη Gauss--Newton} (υποενότητα «\nameref{sec:wls}»). Υπολογίζονται οι
      $\mathbf{h}(\mathbf{x}^{k})$, $\mathbf{H}(\mathbf{x}^{k})$ και $\mathbf{G}(\mathbf{x}^{k})$ και
      επιλύονται οι κανονικές εξισώσεις ως προς τη διόρθωση $\Delta\mathbf{x}^{k}$.
\item \textbf{Έλεγχος σύγκλισης.} Αν $\max_i |\Delta x_i^{k}| \ge \varepsilon$, η κατάσταση
      ενημερώνεται και ο εσωτερικός βρόχος επαναλαμβάνεται. Αν ο αριθμός των επαναλήψεων
      υπερβεί ένα ανώτατο όριο, η εκτέλεση χαρακτηρίζεται μη συγκλίνουσα, κάτι που συνήθως
      υποδηλώνει σφάλμα τοπολογίας, εσφαλμένες παραμέτρους ή χονδροειδές σφάλμα μεγάλου
      μεγέθους.
\item \textbf{Ανίχνευση εσφαλμένων δεδομένων} (υποενότητα «\nameref{sec:baddata}»). Υπολογίζεται η τιμή
      $J(\hat{\mathbf{x}})$ της αντικειμενικής συνάρτησης και εφαρμόζεται ο έλεγχος
      $\chi^2$.
\item \textbf{Αναγνώριση και απόρριψη.} Αν ανιχνευθούν εσφαλμένα δεδομένα, απορρίπτεται η
      μέτρηση με το μέγιστο κανονικοποιημένο υπόλοιπο και η εκτίμηση επαναλαμβάνεται από το
      βήμα 5, με αρχική τιμή την τελευταία λύση. Πριν από την επανεκτίμηση ελέγχεται ότι η
      απόρριψη δεν καθιστά το σύστημα μη παρατηρήσιμο· κάθε απόρριψη μειώνει τον πλεονασμό
      και μπορεί να δημιουργήσει κρίσιμες μετρήσεις (υποενότητα «\nameref{sec:critical}»). Ο εξωτερικός
      βρόχος τερματίζεται όταν ο έλεγχος $\chi^2$ δεν ανιχνεύει πλέον εσφαλμένα δεδομένα.
\item \textbf{Αποτελέσματα.} Η εκτιμώμενη κατάσταση $\hat{\mathbf{x}}$, οι ροές και εγχύσεις
      που υπολογίζονται από αυτήν, η ακρίβεια της εκτίμησης από τα διαγώνια στοιχεία του
      $\mathbf{G}^{-1}$ και ο κατάλογος των απορριφθεισών μετρήσεων διατίθενται στην ανάλυση
      απρόβλεπτων συμβάντων και στις λοιπές εφαρμογές του EMS.
\end{enumerate}

\subsection{Διάνυσμα κατάστασης και μοντέλο μετρήσεων}

Για δίκτυο $N$ ζυγών σε ημιτονοειδή μόνιμη κατάσταση, η κατάσταση λειτουργίας περιγράφεται από
τους φασιθέτες τάσης $\tilde{V}_i = V_i e^{j\theta_i}$ όλων των ζυγών. Σε πολικές συντεταγμένες
ως \emph{διάνυσμα κατάστασης} ορίζεται

\begin{equation}
\mathbf{x} = \left[\theta_2, \ldots, \theta_N,\; V_1, \ldots, V_N\right]^{T}
\label{eq:statevec}
\end{equation}

όπου η γωνία του ζυγού αναφοράς (reference bus) τίθεται $\theta_1 = 0$. Το πλήθος των
μεταβλητών κατάστασης είναι επομένως

\begin{equation}
n = 2N - 1
\label{eq:nstates}
\end{equation}

Όταν διατίθενται μετρήσεις PMU, η αναφορά των γωνιών ορίζεται από τη χρονική βάση UTC, η
γωνία $\theta_1$ εκτιμάται και αυτή και ισχύει $n = 2N$. Από τα $\mathbf{V}$ και
$\boldsymbol{\theta}$ υπολογίζεται κάθε άλλο ηλεκτρικό μέγεθος του δικτύου μέσω των εξισώσεων
ροής φορτίου· το διάνυσμα κατάστασης αποτελεί επομένως το ελάχιστο σύνολο μεταβλητών που
περιγράφει πλήρως τη μόνιμη κατάσταση λειτουργίας.

Το μοντέλο μετρήσεων είναι

\begin{equation}
\mathbf{z} = \mathbf{h}(\mathbf{x}) + \mathbf{e}
\label{eq:measmodel}
\end{equation}

όπου $\mathbf{z} \in \mathbb{R}^{m}$ το διάνυσμα των μετρήσεων, $\mathbf{h}(\cdot)$ το
διάνυσμα των μη γραμμικών συναρτήσεων μετρήσεων και $\mathbf{e}$ το διάνυσμα των σφαλμάτων
μέτρησης. Οι συνήθεις κατηγορίες μετρήσεων είναι:

\begin{itemize}
\item μέτρα τάσεων ζυγών $V_i$·
\item εγχύσεις ενεργού και αέργου ισχύος σε ζυγούς, $P_i$ και $Q_i$·
\item ροές ενεργού και αέργου ισχύος σε κλάδους, $P_{ij}$ και $Q_{ij}$·
\item μέτρα εντάσεων ρεύματος κλάδων $I_{ij}$·
\item φασιθέτες τάσης και ρεύματος από PMU·
\item ψευδομετρήσεις, δηλαδή εκτιμήσεις φορτίου ή παραγωγής με μεγάλη τυπική απόκλιση, και
      εικονικές μετρήσεις (virtual measurements), δηλαδή μηδενικές εγχύσεις σε ζυγούς
      διέλευσης χωρίς παραγωγή και φορτίο, οι οποίες είναι ακριβώς γνωστές.
\end{itemize}

Τα σφάλματα μέτρησης θεωρούνται ανεξάρτητα, κανονικά κατανεμημένα, μηδενικής μέσης τιμής, με
διαγώνιο πίνακα συνδιακύμανσης $\mathbf{R} = \mathrm{diag}(\sigma_1^2, \ldots, \sigma_m^2)$. Η
τυπική απόκλιση κάθε μέτρησης καθορίζεται από την κλάση ακρίβειας των μετασχηματιστών μέτρησης
και των μετατροπέων μέτρησης και από την περιοχή μέτρησης.

Ο \textbf{λόγος πλεονασμού} ορίζεται ως

\begin{equation}
\eta = \frac{m}{n}
\label{eq:redundancy}
\end{equation}

και σε πραγματικά συστήματα κυμαίνεται τυπικά μεταξύ 1,7 και 3. Για παράδειγμα, σε δίκτυο
$N = 100$ ζυγών είναι $n = 199$, οπότε για $\eta = 2$ απαιτούνται περίπου $m = 400$
μετρήσεις. Για $\eta = 1$ το πρόβλημα ανάγεται σε υπολογισμό ροής φορτίου και η ανίχνευση
εσφαλμένων δεδομένων δεν είναι δυνατή.

\subsection{Συναρτήσεις μετρήσεων}

Κάθε κλάδος μεταξύ των ζυγών $i$ και $j$ αναπαρίσταται με το ισοδύναμο κύκλωμα π του
Σχήματος~\ref{fig:pi-model}, με σειριακή αγωγιμότητα $y_{ij} = g_{ij} + j b_{ij} = 1/(r_{ij} +
j x_{ij})$ και εγκάρσια αγωγιμότητα $y_{si} = g_{si} + j b_{si}$ στο άκρο $i$. Για γραμμή
μεταφοράς ισχύει $g_{si} \approx 0$ και η $b_{si}$ ισούται με το ήμισυ της συνολικής
χωρητικής επιδεκτικότητας της γραμμής· οι μετασχηματιστές με λόγο μετασχηματισμού διαφορετικό
του ονομαστικού αναπαρίστανται με ισοδύναμο κύκλωμα π με άνισες εγκάρσιες αγωγιμότητες στα δύο
άκρα. Από τις αγωγιμότητες των κλάδων σχηματίζεται ο πίνακας αγωγιμοτήτων ζυγών
$\mathbf{Y}_{\mathrm{bus}}$, με στοιχεία $G_{ij} + jB_{ij}$.

\begin{figure}[!htbp]
\centering
\includegraphics[width=\linewidth]{figures/pi-model.svg}
\caption{Ισοδύναμο κύκλωμα π κλάδου μεταξύ των ζυγών $i$ και $j$. Η ροή ισχύος
$P_{ij} + jQ_{ij}$ και το ρεύμα $I_{ij}$ ορίζονται στο άκρο $i$.}
\label{fig:pi-model}
\end{figure}

Με $\theta_{ij} = \theta_i - \theta_j$, οι συναρτήσεις μετρήσεων σε πολικές συντεταγμένες
είναι οι ακόλουθες. Για τις \emph{εγχύσεις ισχύος} στον ζυγό $i$, με $\mathcal{N}_i$ το σύνολο
των ζυγών που συνδέονται με τον $i$ συμπεριλαμβανομένου του ίδιου:

\begin{align}
P_i &= V_i \sum_{j \in \mathcal{N}_i} V_j \left( G_{ij}\cos\theta_{ij} + B_{ij}\sin\theta_{ij} \right) \\
Q_i &= V_i \sum_{j \in \mathcal{N}_i} V_j \left( G_{ij}\sin\theta_{ij} - B_{ij}\cos\theta_{ij} \right)
\label{eq:hinj}
\end{align}

Για τις \emph{ροές ισχύος} στον κλάδο $i$--$j$, μετρούμενες στο άκρο $i$:

\begin{align}
P_{ij} &= V_i^2 \left( g_{ij} + g_{si} \right) - V_i V_j \left( g_{ij}\cos\theta_{ij} + b_{ij}\sin\theta_{ij} \right) \\
Q_{ij} &= -V_i^2 \left( b_{ij} + b_{si} \right) - V_i V_j \left( g_{ij}\sin\theta_{ij} - b_{ij}\cos\theta_{ij} \right)
\label{eq:hflow}
\end{align}

Για τα \emph{μέτρα τάσης και ρεύματος}:

\begin{equation}
h_{V_i}(\mathbf{x}) = V_i, \qquad
I_{ij} = \frac{\sqrt{P_{ij}^2 + Q_{ij}^2}}{V_i}
\label{eq:hvi}
\end{equation}

Οι συναρτήσεις \eqref{eq:hinj}--\eqref{eq:hvi} είναι μη γραμμικές: τριγωνομετρικές ως προς τις
γωνίες και δευτέρου βαθμού ως προς τα μέτρα των τάσεων. Οι μετρήσεις μέτρου ρεύματος απαιτούν
ιδιαίτερη μεταχείριση. Σε επίπεδο προφίλ ($V_i = 1$, $\theta_i = 0$) και με αμελητέες
εγκάρσιες αγωγιμότητες ισχύει $P_{ij} = Q_{ij} = 0$, οπότε οι παράγωγοι της $I_{ij}$ δεν
ορίζονται· επιπλέον, η ίδια τιμή $I_{ij}$ αντιστοιχεί σε περισσότερους από έναν συνδυασμούς
$P_{ij}$ και $Q_{ij}$. Για τους λόγους αυτούς οι μετρήσεις μέτρου ρεύματος χρησιμοποιούνται
συνήθως συμπληρωματικά προς τις μετρήσεις ισχύος.

Για τις \emph{μετρήσεις PMU} είναι προτιμότερες οι καρτεσιανές συντεταγμένες,
$\tilde{V}_i = e_i + j f_i$. Ο φασιθέτης τάσης ζυγού αντιστοιχεί τότε απευθείας σε μεταβλητές
κατάστασης και ο φασιθέτης ρεύματος κλάδου είναι γραμμική συνάρτησή τους:

\begin{equation}
\tilde{I}_{ij} = y_{ij}\left(\tilde{V}_i - \tilde{V}_j\right) + y_{si}\,\tilde{V}_i
\label{eq:hpmu}
\end{equation}

\subsection{Εκτίμηση σταθμισμένων ελαχίστων τετραγώνων}
\label{sec:wls}

Για $m > n$ το σύστημα \eqref{eq:measmodel} είναι υπερκαθορισμένο και δεν υπάρχει γενικά
$\mathbf{x}$ που ικανοποιεί όλες τις μετρήσεις. Η εκτίμηση προκύπτει από την ελαχιστοποίηση της
αντικειμενικής συνάρτησης των \textbf{σταθμισμένων ελαχίστων τετραγώνων} (Weighted Least
Squares, WLS):

\begin{equation}
J(\mathbf{x}) = \sum_{i=1}^{m} \frac{\left(z_i - h_i(\mathbf{x})\right)^2}{\sigma_i^2}
             = \left[\mathbf{z}-\mathbf{h}(\mathbf{x})\right]^{T} \mathbf{R}^{-1}
               \left[\mathbf{z}-\mathbf{h}(\mathbf{x})\right]
\label{eq:wls}
\end{equation}

Ο συντελεστής στάθμισης $1/\sigma_i^2$ αποδίδει μεγαλύτερο βάρος στις ακριβέστερες μετρήσεις.
Κάθε όρος του αθροίσματος είναι το τετράγωνο του υπολοίπου εκφρασμένου σε μονάδες τυπικής
απόκλισης, ώστε μετρήσεις διαφορετικών μεγεθών και κλιμάκων να είναι συγκρίσιμες. Για
κανονικά κατανεμημένα σφάλματα, η ελαχιστοποίηση της \eqref{eq:wls} ταυτίζεται με την
εκτίμηση μέγιστης πιθανοφάνειας.

Στο ελάχιστο μηδενίζεται η κλίση της αντικειμενικής συνάρτησης:

\begin{equation}
\mathbf{g}(\mathbf{x}) = \frac{\partial J}{\partial \mathbf{x}}
  = -2\,\mathbf{H}^{T}(\mathbf{x})\,\mathbf{R}^{-1}\left[\mathbf{z} - \mathbf{h}(\mathbf{x})\right]
  = \mathbf{0}, \qquad
\mathbf{H}(\mathbf{x}) = \frac{\partial \mathbf{h}}{\partial \mathbf{x}}
\label{eq:foc}
\end{equation}

όπου $\mathbf{H}(\mathbf{x}) \in \mathbb{R}^{m \times n}$ ο ιακωβιανός πίνακας των συναρτήσεων
μετρήσεων. Η \eqref{eq:foc} είναι σύστημα $n$ μη γραμμικών εξισώσεων. Η μέθοδος Newton θα
απαιτούσε τον εσσιανό πίνακα της $J$, ο οποίος περιέχει τις δεύτερες παραγώγους της
$\mathbf{h}$. Η μέθοδος Gauss--Newton προκύπτει από τη γραμμικοποίηση
$\mathbf{h}(\mathbf{x}^{k} + \Delta\mathbf{x}) \approx \mathbf{h}(\mathbf{x}^{k}) +
\mathbf{H}(\mathbf{x}^{k})\,\Delta\mathbf{x}$ και την ελαχιστοποίηση της $J$ ως προς
$\Delta\mathbf{x}$, δηλαδή με παράλειψη των όρων δεύτερης τάξης. Σε κάθε επανάληψη $k$
επιλύονται οι \emph{κανονικές εξισώσεις}

\begin{align}
\mathbf{G}(\mathbf{x}^{k}) \, \Delta \mathbf{x}^{k}
  &= \mathbf{H}^{T}(\mathbf{x}^{k}) \, \mathbf{R}^{-1}
     \left[\mathbf{z} - \mathbf{h}(\mathbf{x}^{k})\right] \\
\mathbf{G}(\mathbf{x}^{k})
  &= \mathbf{H}^{T}(\mathbf{x}^{k}) \, \mathbf{R}^{-1} \mathbf{H}(\mathbf{x}^{k})
\label{eq:gain}
\end{align}

Η παράλειψη των δεύτερων παραγώγων είναι ακριβής για γραμμικό μοντέλο μετρήσεων και
επιτρεπτή όταν τα υπόλοιπα στη λύση είναι μικρά· στην περίπτωση αυτή η σύγκλιση είναι σχεδόν
τετραγωνική. Ο αλγόριθμος περιλαμβάνει τα ακόλουθα βήματα:

\begin{enumerate}
\item Αρχικοποίηση: $\mathbf{x}^{0}$ ίσο με το επίπεδο προφίλ ($V_i = 1$, $\theta_i = 0$) ή
      με τη λύση της προηγούμενης εκτέλεσης· $k = 0$.
\item Υπολογισμός των $\mathbf{h}(\mathbf{x}^{k})$, $\mathbf{H}(\mathbf{x}^{k})$ και
      $\mathbf{G}(\mathbf{x}^{k})$.
\item Επίλυση των κανονικών εξισώσεων \eqref{eq:gain} ως προς $\Delta\mathbf{x}^{k}$.
\item Αν $\max_i |\Delta x_i^{k}| < \varepsilon$, τερματισμός· διαφορετικά
      $\mathbf{x}^{k+1} = \mathbf{x}^{k} + \Delta\mathbf{x}^{k}$, $k \leftarrow k + 1$ και
      επιστροφή στο βήμα 2.
\end{enumerate}

Η ανοχή $\varepsilon$ είναι τυπικά της τάξης του $10^{-4}$ (p.u. για τα μέτρα τάσεων, rad για
τις γωνίες). Από επίπεδο προφίλ η σύγκλιση επιτυγχάνεται συνήθως σε τρεις έως έξι επαναλήψεις
και από τη λύση της προηγούμενης εκτέλεσης σε λιγότερες.

Ο πίνακας $\mathbf{G}$ ονομάζεται \emph{πίνακας κέρδους} (gain matrix). Είναι συμμετρικός,
θετικά ορισμένος όταν το σύστημα είναι παρατηρήσιμο, και \emph{αραιός}, επειδή κάθε μέτρηση
συνδέεται μόνο με τις μεταβλητές κατάστασης γειτονικών ζυγών. Οι κανονικές εξισώσεις
επιλύονται με αραιή παραγοντοποίηση Cholesky, της οποίας το υπολογιστικό κόστος αυξάνει
σχεδόν γραμμικά με το μέγεθος του δικτύου, έναντι $O(n^3)$ για πυκνή παραγοντοποίηση. Η
αραιότητα καθιστά επομένως δυνατή την εκτέλεση σε πραγματικό χρόνο για δίκτυα χιλιάδων
ζυγών. Στη λύση, ο $\mathbf{G}^{-1}(\hat{\mathbf{x}})$ αποτελεί προσέγγιση του πίνακα
συνδιακύμανσης του σφάλματος εκτίμησης $\hat{\mathbf{x}} - \mathbf{x}$· τα διαγώνια στοιχεία
του δίνουν την αναμενόμενη ακρίβεια κάθε μεταβλητής κατάστασης.

\subsection{Δομή του ιακωβιανού πίνακα}

Με διάταξη των μεταβλητών κατάστασης κατά \eqref{eq:statevec} και των μετρήσεων κατά
κατηγορία, ο ιακωβιανός πίνακας έχει τη δομή μπλοκ

\begin{equation}
\mathbf{H} =
\begin{bmatrix}
\partial \mathbf{P}_{\mathrm{inj}} / \partial \boldsymbol{\theta} & \partial \mathbf{P}_{\mathrm{inj}} / \partial \mathbf{V} \\
\partial \mathbf{P}_{\mathrm{flow}} / \partial \boldsymbol{\theta} & \partial \mathbf{P}_{\mathrm{flow}} / \partial \mathbf{V} \\
\partial \mathbf{Q}_{\mathrm{inj}} / \partial \boldsymbol{\theta} & \partial \mathbf{Q}_{\mathrm{inj}} / \partial \mathbf{V} \\
\partial \mathbf{Q}_{\mathrm{flow}} / \partial \boldsymbol{\theta} & \partial \mathbf{Q}_{\mathrm{flow}} / \partial \mathbf{V} \\
\mathbf{0} & \partial \mathbf{V}_{\mathrm{meas}} / \partial \mathbf{V}
\end{bmatrix}
\label{eq:jacblocks}
\end{equation}

όπου η τελευταία γραμμή έχει μοναδιαία στοιχεία στις στήλες των μετρούμενων τάσεων. Από τις
\eqref{eq:hinj}, τα στοιχεία των γραμμών εγχύσεων είναι, για $j \neq i$:

\begin{align}
\frac{\partial P_i}{\partial \theta_i} &= -Q_i - B_{ii} V_i^2 \\
\frac{\partial P_i}{\partial \theta_j} &= V_i V_j \left( G_{ij}\sin\theta_{ij} - B_{ij}\cos\theta_{ij} \right) \\
\frac{\partial P_i}{\partial V_i} &= \frac{P_i}{V_i} + G_{ii} V_i \\
\frac{\partial P_i}{\partial V_j} &= V_i \left( G_{ij}\cos\theta_{ij} + B_{ij}\sin\theta_{ij} \right) \\
\frac{\partial Q_i}{\partial \theta_i} &= P_i - G_{ii} V_i^2 \\
\frac{\partial Q_i}{\partial \theta_j} &= -V_i V_j \left( G_{ij}\cos\theta_{ij} + B_{ij}\sin\theta_{ij} \right) \\
\frac{\partial Q_i}{\partial V_i} &= \frac{Q_i}{V_i} - B_{ii} V_i \\
\frac{\partial Q_i}{\partial V_j} &= V_i \left( G_{ij}\sin\theta_{ij} - B_{ij}\cos\theta_{ij} \right)
\label{eq:jacinj}
\end{align}

Από τις \eqref{eq:hflow}, τα στοιχεία των γραμμών ροών είναι:

\begin{align}
\frac{\partial P_{ij}}{\partial \theta_i} &= -\frac{\partial P_{ij}}{\partial \theta_j}
  = V_i V_j \left( g_{ij}\sin\theta_{ij} - b_{ij}\cos\theta_{ij} \right) \\
\frac{\partial P_{ij}}{\partial V_i} &= -V_j \left( g_{ij}\cos\theta_{ij} + b_{ij}\sin\theta_{ij} \right) + 2\left( g_{ij} + g_{si} \right) V_i \\
\frac{\partial P_{ij}}{\partial V_j} &= -V_i \left( g_{ij}\cos\theta_{ij} + b_{ij}\sin\theta_{ij} \right) \\
\frac{\partial Q_{ij}}{\partial \theta_i} &= -\frac{\partial Q_{ij}}{\partial \theta_j}
  = -V_i V_j \left( g_{ij}\cos\theta_{ij} + b_{ij}\sin\theta_{ij} \right) \\
\frac{\partial Q_{ij}}{\partial V_i} &= -V_j \left( g_{ij}\sin\theta_{ij} - b_{ij}\cos\theta_{ij} \right) - 2\left( b_{ij} + b_{si} \right) V_i \\
\frac{\partial Q_{ij}}{\partial V_j} &= -V_i \left( g_{ij}\sin\theta_{ij} - b_{ij}\cos\theta_{ij} \right)
\label{eq:jacflow}
\end{align}

Ο ιακωβιανός είναι αραιός: η γραμμή μιας μέτρησης έγχυσης στον ζυγό $i$ έχει μη μηδενικά
στοιχεία μόνο στις στήλες του ζυγού $i$ και των γειτονικών του ζυγών, ενώ η γραμμή μιας
μέτρησης ροής μόνο στις στήλες των ζυγών $i$ και $j$. Για γραμμές μεταφοράς με
$r_{ij} \ll x_{ij}$ και μικρές διαφορές γωνιών ισχύει $|g_{ij}| \ll |b_{ij}|$,
$\cos\theta_{ij} \approx 1$ και $\sin\theta_{ij} \approx 0$. Τα στοιχεία των μπλοκ
$\partial\mathbf{P}/\partial\boldsymbol{\theta}$ και $\partial\mathbf{Q}/\partial\mathbf{V}$
είναι τότε της τάξης του $|b_{ij}|$, ενώ εκείνα των $\partial\mathbf{P}/\partial\mathbf{V}$ και
$\partial\mathbf{Q}/\partial\boldsymbol{\theta}$ είναι σημαντικά μικρότερα: η ενεργός ισχύς
συνδέεται ισχυρά με τις γωνίες και η άεργος ισχύς με τα μέτρα των τάσεων.

\subsection{Παραλλαγές της εκτίμησης WLS}

\textbf{Γραμμικό μοντέλο ΣΡ.} Με τις παραδοχές $V_i \approx 1$, $r_{ij} \ll x_{ij}$ (οπότε
$g_{ij} \approx 0$ και $b_{ij} \approx -1/x_{ij}$), $\sin\theta_{ij} \approx \theta_{ij}$,
$\cos\theta_{ij} \approx 1$ και αμελητέες εγκάρσιες αγωγιμότητες, η \eqref{eq:hflow} για την
ενεργό ισχύ ανάγεται στη γραμμική σχέση

\begin{equation}
P_{ij} = \frac{\theta_i - \theta_j}{x_{ij}}
\label{eq:dcflow}
\end{equation}

Οι μετρήσεις ενεργού ισχύος εκφράζονται τότε ως $\mathbf{z}_P = \mathbf{H}_{\mathrm{DC}}\,
\boldsymbol{\theta} + \mathbf{e}_P$ με σταθερό ιακωβιανό, και οι γωνίες υπολογίζονται χωρίς
επαναλήψεις ως $\hat{\boldsymbol{\theta}} = (\mathbf{H}_{\mathrm{DC}}^{T}\mathbf{R}_P^{-1}
\mathbf{H}_{\mathrm{DC}})^{-1}\mathbf{H}_{\mathrm{DC}}^{T}\mathbf{R}_P^{-1}\mathbf{z}_P$. Το
μοντέλο ΣΡ δεν εκτιμά μέτρα τάσεων και άεργο ισχύ, και η ακρίβειά του μειώνεται με την αύξηση
του λόγου $r/x$ και της φόρτισης. Χρησιμοποιείται στην ανάλυση παρατηρησιμότητας και για
τον υπολογισμό αρχικής εκτίμησης.

\textbf{Αποζευγμένος εκτιμητής.} Με βάση την ασθενή σύζευξη μεταξύ ενεργού ισχύος και μέτρων
τάσεων και μεταξύ αέργου ισχύος και γωνιών, τα μη διαγώνια μπλοκ του $\mathbf{G}$ παραλείπονται
και οι κανονικές εξισώσεις διασπώνται σε δύο μικρότερα συστήματα: ένα για τις γωνίες, από
τις μετρήσεις ενεργού ισχύος, και ένα για τα μέτρα των τάσεων, από τις μετρήσεις αέργου
ισχύος και τάσης. Στον ταχύ αποζευγμένο εκτιμητή (fast decoupled estimator) τα μπλοκ
$\mathbf{G}_{P\theta}$ και $\mathbf{G}_{QV}$ υπολογίζονται μία φορά σε επίπεδο προφίλ και
παραμένουν σταθερά, οπότε παραγοντοποιούνται μόνο μία φορά. Αν οι προσεγγίσεις περιορίζονται
στον πίνακα κέρδους και το δεξί μέλος υπολογίζεται με τον ακριβή ιακωβιανό, η λύση ταυτίζεται
με εκείνη του πλήρους εκτιμητή, με περισσότερες αλλά φθηνότερες επαναλήψεις· προσεγγίσεις και
στο δεξί μέλος οδηγούν σε ελαφρώς διαφορετική λύση.

\textbf{Εικονικές μετρήσεις και περιορισμοί ισότητας.} Οι μηδενικές εγχύσεις των ζυγών
διέλευσης μπορούν να εισαχθούν ως μετρήσεις πολύ μικρής τυπικής απόκλισης. Τα αντίστοιχα βάρη
$1/\sigma_i^2$ είναι όμως κατά τάξεις μεγέθους μεγαλύτερα από τα βάρη των πραγματικών
μετρήσεων και αυξάνουν δραστικά τον δείκτη κατάστασης του $\mathbf{G}$. Εναλλακτικά, οι
εικονικές μετρήσεις διατυπώνονται ως περιορισμοί ισότητας $\mathbf{c}(\mathbf{x}) = \mathbf{0}$
και σε κάθε επανάληψη επιλύεται, με πολλαπλασιαστές Lagrange $\boldsymbol{\lambda}$, το σύστημα

\begin{equation}
\begin{bmatrix} \mathbf{G} & \mathbf{C}^{T} \\ \mathbf{C} & \mathbf{0} \end{bmatrix}
\begin{bmatrix} \Delta\mathbf{x}^{k} \\ \boldsymbol{\lambda} \end{bmatrix}
=
\begin{bmatrix} \mathbf{H}^{T}\mathbf{R}^{-1}\left[\mathbf{z} - \mathbf{h}(\mathbf{x}^{k})\right] \\ -\mathbf{c}(\mathbf{x}^{k}) \end{bmatrix}
\label{eq:eqconstr}
\end{equation}

όπου $\mathbf{C} = \partial\mathbf{c}/\partial\mathbf{x}$ και όλοι οι πίνακες υπολογίζονται στο
$\mathbf{x}^{k}$. Ο πίνακας συντελεστών είναι συμμετρικός αλλά όχι θετικά ορισμένος, οπότε η
παραγοντοποίησή του απαιτεί κατάλληλη επιλογή οδηγών (pivoting).

\textbf{Αριθμητική ευστάθεια.} Με $\tilde{\mathbf{H}} = \mathbf{R}^{-1/2}\mathbf{H}$ τον
σταθμισμένο ιακωβιανό ισχύει $\mathbf{G} = \tilde{\mathbf{H}}^{T}\tilde{\mathbf{H}}$ και
$\kappa(\mathbf{G}) = \kappa(\tilde{\mathbf{H}})^2$, όπου $\kappa$ ο δείκτης κατάστασης. Ο
σχηματισμός του $\mathbf{G}$ επιδεινώνει επομένως την αριθμητική συμπεριφορά, ιδίως όταν
συνυπάρχουν μετρήσεις πολύ διαφορετικής ακρίβειας ή κλάδοι με πολύ διαφορετικές αντιδράσεις.
Η ορθογώνια παραγοντοποίηση $\tilde{\mathbf{H}} = \mathbf{Q}\mathbf{U}$, με ορθογώνιο πίνακα
$\mathbf{Q}$ και άνω τριγωνικό $\mathbf{U}$ (για παράδειγμα με περιστροφές Givens), επιλύει το
πρόβλημα ελαχίστων τετραγώνων χωρίς σχηματισμό του $\mathbf{G}$, αφού
$\mathbf{U}^{T}\mathbf{U} = \mathbf{G}$. Η μέθοδος του επαυξημένου πίνακα (Hachtel) εισάγει τα
υπόλοιπα ως πρόσθετους αγνώστους και επιλύει ένα μεγαλύτερο αλλά αραιό και καλύτερα
ρυθμισμένο σύστημα, στο οποίο οι εικονικές μετρήσεις ενσωματώνονται ως περιορισμοί.

\textbf{Εύρωστοι εκτιμητές.} Η εκτίμηση WLS είναι ευαίσθητη στα χονδροειδή σφάλματα, επειδή
το τετραγωνικό κριτήριο αποδίδει μεγάλο βάρος στα μεγάλα υπόλοιπα. Ο εκτιμητής σταθμισμένων
ελαχίστων απολύτων τιμών (Weighted Least Absolute Value, WLAV) ελαχιστοποιεί το
$\sum_{i} |z_i - h_i(\mathbf{x})|/\sigma_i$ και επιλύεται με γραμμικό προγραμματισμό. Η λύση
του ικανοποιεί ακριβώς $n$ μετρήσεις και απορρίπτει αυτόματα τα μεμονωμένα εσφαλμένα
δεδομένα, εκτός αν αυτά αφορούν σημεία μόχλευσης (leverage points), δηλαδή μετρήσεις με
ασυνήθιστα μεγάλα στοιχεία στον ιακωβιανό, όπως εγχύσεις σε ζυγούς με πολλούς κλάδους ή
κλάδους πολύ μικρής αντίδρασης.

\subsection{Παρατηρησιμότητα}
\label{sec:observability}

Η επίλυση των κανονικών εξισώσεων \eqref{eq:gain} προϋποθέτει ότι ο $\mathbf{G}$ είναι
αντιστρέψιμος. Η συνθήκη διατυπώνεται με την έννοια της τάξης του ιακωβιανού πίνακα.

\textbf{Τάξη πίνακα.} Η \emph{τάξη} (rank) ενός πίνακα $\mathbf{A} \in \mathbb{R}^{m \times n}$
είναι το μέγιστο πλήθος γραμμικώς ανεξάρτητων στηλών του, το οποίο ισούται με το μέγιστο
πλήθος γραμμικώς ανεξάρτητων γραμμών του· ισοδύναμα, είναι η διάσταση του διανυσματικού χώρου
που παράγουν οι στήλες του. Ισχύει $\mathrm{rank}(\mathbf{A}) \le \min(m, n)$. Ο $\mathbf{A}$
έχει \emph{πλήρη τάξη στηλών} όταν $\mathrm{rank}(\mathbf{A}) = n$, δηλαδή όταν η εξίσωση
$\mathbf{A}\mathbf{v} = \mathbf{0}$ έχει μόνο τη λύση $\mathbf{v} = \mathbf{0}$.

Για τον ιακωβιανό, η συνθήκη $\mathrm{rank}(\mathbf{H}) = n$ σημαίνει ότι καμία μη μηδενική
μεταβολή $\Delta\mathbf{x}$ της κατάστασης δεν αφήνει αμετάβλητες όλες τις γραμμικοποιημένες
μετρήσεις. Επειδή
$\Delta\mathbf{x}^{T}\mathbf{G}\,\Delta\mathbf{x} = \|\mathbf{R}^{-1/2}\mathbf{H}\,\Delta\mathbf{x}\|^{2}$,
ο $\mathbf{G}$ είναι θετικά ορισμένος, και επομένως αντιστρέψιμος, αν και μόνο αν ο
$\mathbf{H}$ έχει πλήρη τάξη στηλών. Η συνθήκη $m \ge n$ είναι αναγκαία αλλά όχι ικανή. Στο
δίκτυο τριών ζυγών των παραδειγμάτων με μοντέλο ΣΡ, για παράδειγμα, οι μετρήσεις $P_{12}$ και
$P_{21}$ της ροής του ίδιου κλάδου στα δύο άκρα του δίνουν $m = n = 2$, αλλά οι γραμμές
$[-1/x_{12},\; 0]$ και $[1/x_{12},\; 0]$ του $\mathbf{H}$ είναι γραμμικώς εξαρτημένες: ισχύει
$\mathrm{rank}(\mathbf{H}) = 1$ και η γωνία $\theta_3$ δεν προσδιορίζεται. Η τάξη εξαρτάται
επομένως από τη θέση και το είδος των μετρήσεων και όχι μόνο από το πλήθος τους: μετρήσεις
συγκεντρωμένες σε τμήμα του δικτύου δεν επιτρέπουν τον προσδιορισμό της κατάστασης στο
υπόλοιπο δίκτυο.

Το σύστημα χαρακτηρίζεται \textbf{παρατηρήσιμο} όταν οι διαθέσιμες μετρήσεις επιτρέπουν τον
μονοσήμαντο προσδιορισμό ολόκληρου του διανύσματος κατάστασης. Η παρατηρησιμότητα ελέγχεται
με δύο μεθόδους:

\begin{itemize}
\item \textbf{Τοπολογική}: αναζητείται επικαλύπτον δένδρο (spanning tree) του γράφου του
      δικτύου, κάθε κλάδος του οποίου αντιστοιχίζεται σε μέτρηση. Βασίζεται στη θεωρία
      γράφων και δεν επηρεάζεται από σφάλματα στρογγυλοποίησης.
\item \textbf{Αριθμητική}: ελέγχεται η τάξη του $\mathbf{H}$ μέσω της τριγωνικής
      παραγοντοποίησης του $\mathbf{G}$, με εντοπισμό μηδενικών οδηγών (pivots). Επειδή τα
      σφάλματα στρογγυλοποίησης σπάνια δίνουν ακριβώς μηδέν, ένας οδηγός θεωρείται μηδενικός
      όταν είναι μικρότερος από ένα κατώφλι. Η μέθοδος είναι γενικότερη και αντιμετωπίζει και
      τις οριακές περιπτώσεις.
\end{itemize}

Και οι δύο μέθοδοι εφαρμόζονται συνήθως στο αποζευγμένο γραμμικό μοντέλο, χωριστά για τα
ζεύγη $P$--$\theta$ και $Q$--$V$. Σε μη παρατηρήσιμο σύστημα προσδιορίζονται οι
\emph{παρατηρήσιμες νησίδες} (observable islands) και η κατάσταση εκτιμάται εντός αυτών. Η
παρατηρησιμότητα αποκαθίσταται με \emph{ψευδομετρήσεις} (pseudo-measurements), δηλαδή
εκτιμήσεις φορτίου από ιστορικά δεδομένα ή προβλέψεις, στις οποίες αποδίδεται μεγάλη τυπική
απόκλιση ώστε η επίδρασή τους στην εκτίμηση να είναι μικρή.

\subsection{Ανίχνευση και αναγνώριση εσφαλμένων μετρήσεων}
\label{sec:baddata}

Η ανίχνευση εσφαλμένων δεδομένων (bad data) βασίζεται στα \emph{υπόλοιπα} (residuals)
$\mathbf{r} = \mathbf{z}-\mathbf{h}(\hat{\mathbf{x}})$ της εκτίμησης. Με γραμμικοποίηση γύρω από
τη λύση, τα υπόλοιπα συνδέονται με τα σφάλματα μέτρησης μέσω του πίνακα ευαισθησίας
υπολοίπων, $\mathbf{r} \approx \mathbf{S}\mathbf{e}$ με
$\mathbf{S} = \mathbf{I} - \mathbf{H}\mathbf{G}^{-1}\mathbf{H}^{T}\mathbf{R}^{-1}$. Η κλασική
διαδικασία περιλαμβάνει δύο στάδια.

\textbf{Στάδιο 1 --- Ανίχνευση.} Όταν δεν υπάρχουν εσφαλμένα δεδομένα, η τιμή
$J(\hat{\mathbf{x}})$ ακολουθεί προσεγγιστικά κατανομή $\chi^2$ με $m-n$ βαθμούς ελευθερίας
και μέση τιμή $m-n$· για το προηγούμενο παράδειγμα ($m = 400$, $n = 199$) η αναμενόμενη τιμή
είναι 201. Αν $J(\hat{\mathbf{x}}) > \chi^2_{m-n,\,p}$, όπου $p$ η πιθανότητα του ελέγχου
(συνήθως 95\%), θεωρείται ότι υπάρχουν εσφαλμένα δεδομένα. Ο έλεγχος $\chi^2$ ανιχνεύει την
ύπαρξη εσφαλμένων δεδομένων, αλλά δεν τα εντοπίζει.

\textbf{Στάδιο 2 --- Αναγνώριση.} Τα υπόλοιπα έχουν διαφορετικές διασπορές και δεν
συγκρίνονται απευθείας. Κανονικοποιούνται ως

\begin{equation}
r_i^{N} = \frac{|r_i|}{\sqrt{\Omega_{ii}}}, \qquad
\boldsymbol{\Omega} = \mathbf{S}\mathbf{R} = \mathbf{R} - \mathbf{H}\mathbf{G}^{-1}\mathbf{H}^{T}
\label{eq:normres}
\end{equation}

όπου $\boldsymbol{\Omega}$ ο πίνακας συνδιακύμανσης των υπολοίπων. Κατά το κριτήριο του
μεγίστου κανονικοποιημένου υπολοίπου (Largest Normalized Residual, LNR), αν το μέγιστο
$r_i^{N}$ υπερβαίνει ένα κατώφλι, συνήθως 3, η αντίστοιχη μέτρηση απορρίπτεται και η
εκτίμηση επαναλαμβάνεται. Η διαδικασία συνεχίζεται μέχρι να μην υπάρχει κανονικοποιημένο
υπόλοιπο μεγαλύτερο από το κατώφλι. Το κριτήριο LNR είναι αξιόπιστο για μεμονωμένα
εσφαλμένα δεδομένα, αλλά μπορεί να αστοχήσει όταν υπάρχουν πολλαπλά εσφαλμένα δεδομένα σε
μετρήσεις με ισχυρά συσχετισμένα υπόλοιπα.

\subsection{Κρίσιμες μετρήσεις}
\label{sec:critical}

Η ανίχνευση εσφαλμένων δεδομένων μέσω των υπολοίπων προϋποθέτει πλεονασμό όχι μόνο συνολικά
αλλά και τοπικά, στην περιοχή κάθε μέτρησης.

Μια μέτρηση χαρακτηρίζεται \emph{κρίσιμη} (critical measurement) όταν η αφαίρεσή της καθιστά το
σύστημα μη παρατηρήσιμο, δηλαδή όταν μειώνει την τάξη του $\mathbf{H}$ κάτω από $n$. Για
κρίσιμη μέτρηση ισχύει $\Omega_{ii} = 0$ και το υπόλοιπό της είναι πάντοτε μηδενικό, επειδή
καμία άλλη μέτρηση δεν παρέχει πληροφορία για το αντίστοιχο μέγεθος: η εκτίμηση προσαρμόζεται
ακριβώς στην τιμή της, όποια κι αν είναι αυτή. Σφάλμα σε κρίσιμη μέτρηση είναι επομένως μη
ανιχνεύσιμο, ανεξάρτητα από το μέγεθός του, και μεταφέρεται αυτούσιο στην εκτιμώμενη
κατάσταση.

Αντίστοιχα, δύο μετρήσεις αποτελούν \emph{κρίσιμο ζεύγος} (critical pair) όταν καμία δεν είναι
κρίσιμη, αλλά η ταυτόχρονη αφαίρεσή τους καθιστά το σύστημα μη παρατηρήσιμο. Τα υπόλοιπα των
δύο μετρήσεων είναι πλήρως συσχετισμένα και τα κανονικοποιημένα υπόλοιπά τους ίσα: σφάλμα σε
μία από αυτές ανιχνεύεται από τον έλεγχο $\chi^2$, αλλά δεν αναγνωρίζεται με το κριτήριο LNR,
επειδή δεν είναι δυνατόν να διακριθεί ποια από τις δύο είναι εσφαλμένη.

Η αποφυγή κρίσιμων μετρήσεων και κρίσιμων ζευγών αποτελεί κριτήριο σχεδιασμού του συστήματος
μετρήσεων, ιδίως επειδή η απώλεια μετρήσεων λόγω βλάβης ή η απόρριψή τους ως εσφαλμένων
μπορεί να μετατρέψει πλεονάζουσες μετρήσεις σε κρίσιμες. Η περίπτωση αναλύεται αριθμητικά
στο πρώτο από τα Παραδείγματα του κεφαλαίου, στην ενότητα «Κρίσιμες μετρήσεις».

\subsection{Ενσωμάτωση μετρήσεων PMU}

Οι PMU μετρούν απευθείας φασιθέτες τάσης και ρεύματος, οι οποίοι σε καρτεσιανές συντεταγμένες
είναι γραμμικές συναρτήσεις της κατάστασης σύμφωνα με την \eqref{eq:hpmu}. Αν το σύστημα είναι
πλήρως παρατηρήσιμο μόνο από PMU, το μοντέλο \eqref{eq:measmodel} γίνεται γραμμικό,
$\mathbf{z} = \mathbf{H}\mathbf{x} + \mathbf{e}$, και η εκτίμηση υπολογίζεται χωρίς
επαναλήψεις:

\begin{equation}
\hat{\mathbf{x}} = \left(\mathbf{H}^{T}\mathbf{R}^{-1}\mathbf{H}\right)^{-1}
                    \mathbf{H}^{T}\mathbf{R}^{-1}\mathbf{z}
\label{eq:linse}
\end{equation}

Η γραμμική εκτίμηση δεν παρουσιάζει προβλήματα σύγκλισης ούτε εξάρτηση από την αρχική
εκτίμηση. Ο $\mathbf{H}$ είναι σταθερός για δεδομένη τοπολογία, οπότε ο πίνακας κέρδους
παραγοντοποιείται μία φορά, και η εκτίμηση μπορεί να εκτελείται με τον ρυθμό αναφοράς των
PMU, δηλαδή δεκάδες φορές ανά δευτερόλεπτο, αντί για μία φορά ανά ένα έως λίγα λεπτά. Επειδή
οι PMU αναφέρουν συνήθως μέτρο και γωνία, η μετατροπή σε καρτεσιανές συντεταγμένες συσχετίζει
τα σφάλματα των δύο συνιστωσών, οπότε ο αντίστοιχος πίνακας συνδιακύμανσης δεν είναι
διαγώνιος.

Επειδή η πλήρης κάλυψη με PMU δεν είναι συνήθως οικονομικά εφικτή, χρησιμοποιούνται
\textbf{υβριδικοί εκτιμητές} που συνδυάζουν μετρήσεις SCADA και PMU. Η κύρια δυσκολία τους
είναι ο συνδυασμός μετρήσεων με διαφορετική ακρίβεια και διαφορετικό ρυθμό ανανέωσης,
ενδεικτικά 2 έως 10 s για το SCADA έναντι 20 ms για τις PMU.

\section{Ευέλικτα συστήματα μεταφοράς εναλλασσόμενου ρεύματος (FACTS)}

Πέρα από τη ρύθμιση της παραγωγής μέσω των αγορών και των εφεδρειών συχνότητας, ο ΔΣΜ
ελέγχει άμεσα τις ροές ισχύος και τις τάσεις του δικτύου μεταφοράς. Οι διατάξεις
\textbf{FACTS} (Flexible AC Transmission Systems) είναι διατάξεις ηλεκτρονικών ισχύος που
επιτρέπουν τον ταχύ και συνεχή έλεγχο των μεγεθών που καθορίζουν τη ροή ισχύος σε δίκτυο
εναλλασσόμενου ρεύματος.

Σύμφωνα με την \eqref{eq:powerflow}, η ροή ενεργού ισχύος εξαρτάται από τα μέτρα των τάσεων
$V_1, V_2$, την αντίδραση σειράς $X$ και τη διαφορά γωνιών $\delta_1 - \delta_2$. Σε
συμβατικό δίκτυο τα μεγέθη αυτά δεν ελέγχονται άμεσα, οπότε η κατανομή των ροών καθορίζεται
από τις σύνθετες αντιστάσεις των κλάδων και όχι από τις απαιτήσεις λειτουργίας. Οι διατάξεις
FACTS καθιστούν ελεγχόμενο ένα ή περισσότερα από τα μεγέθη αυτά και κατατάσσονται ανάλογα με
το μέγεθος που επηρεάζουν (Σχήμα~\ref{fig:facts}):

\begin{itemize}
\item \textbf{Παράλληλης σύνδεσης} (έλεγχος του $V$): ο στατός αντισταθμιστής αέργου ισχύος
      (Static Var Compensator, SVC) και ο στατός σύγχρονος αντισταθμιστής (Static
      Synchronous Compensator, STATCOM). Ρυθμίζουν την τάση του ζυγού μέσω έγχυσης ή
      απορρόφησης αέργου ισχύος.
\item \textbf{Σειριακής σύνδεσης} (έλεγχος του $X$): ο πυκνωτής σειράς ελεγχόμενος με
      θυρίστορ (Thyristor Controlled Series Capacitor, TCSC) και ο στατός σύγχρονος
      αντισταθμιστής σειράς (Static Synchronous Series Compensator, SSSC). Μεταβάλλουν την
      ισοδύναμη αντίδραση σειράς της γραμμής, ανακατανέμουν τις ροές ισχύος και αυξάνουν τη
      μεταφορική ικανότητα.
\item \textbf{Συνδυασμένης σύνδεσης}, σε σειρά και παράλληλα (έλεγχος και της γωνίας): ο
      ενοποιημένος ελεγκτής ροής ισχύος (Unified Power Flow Controller, UPFC) ελέγχει
      ταυτόχρονα και ανεξάρτητα την τάση και τις ροές ενεργού και αέργου ισχύος. Παρέχει τη
      μεγαλύτερη ευελιξία ελέγχου, με το υψηλότερο κόστος.
\end{itemize}

\begin{figure}[!htbp]
\centering
\includegraphics[width=\linewidth]{figures/facts.svg}
\caption{Οι τρεις κατηγορίες διατάξεων FACTS σε γραμμή δύο ζυγών. Κάθε κατηγορία επεμβαίνει
σε διαφορετικό μέγεθος της σχέσης μεταφοράς ισχύος, το οποίο σημειώνεται με έντονο χρώμα: η
παράλληλη στην τάση, η σειριακή στην αντίδραση, η συνδυασμένη και στη γωνία.}
\label{fig:facts}
\end{figure}

\subsection{Τοπολογίες διατάξεων}

Ανάλογα με την αρχή λειτουργίας τους, οι διατάξεις FACTS διακρίνονται σε διατάξεις
\emph{μεταβλητής σύνθετης αντίστασης} (εμπέδησης), με θυρίστορ, και σε διατάξεις
\emph{ελεγχόμενης πηγής τάσης}, με μετατροπείς πηγής τάσης. Τα μονογραμμικά διαγράμματα
δίνονται στο Σχήμα~\ref{fig:facts-devices}.

\begin{itemize}
\item Ο \textbf{SVC} αποτελείται από πηνίο ελεγχόμενο με θυρίστορ (Thyristor Controlled
      Reactor, TCR) και πυκνωτές μεταγόμενους με θυρίστορ (Thyristor Switched Capacitor,
      TSC), σε παράλληλη σύνδεση. Η γωνία έναυσης του TCR ρυθμίζει συνεχώς την ισοδύναμη
      επαγωγική επιδεκτικότητα, ενώ οι TSC προσθέτουν χωρητική επιδεκτικότητα σε διακριτά
      βήματα. Η διάταξη συμπεριφέρεται ως \emph{μεταβλητή επιδεκτικότητα} $B$ (susceptance).
\item Ο \textbf{STATCOM} είναι μετατροπέας πηγής τάσης (Voltage Source Converter, VSC) με
      πυκνωτή στην πλευρά συνεχούς ρεύματος (ΣΡ). Παράγει εναλλασσόμενη τάση $V_\mu$
      ελεγχόμενου μέτρου, η οποία εφαρμόζεται στο δίκτυο μέσω της αντίδρασης σύζευξης: για
      $V_\mu > V$ ο STATCOM εγχέει άεργο ισχύ, ενώ για $V_\mu < V$ απορροφά.
\item Ο \textbf{TCSC} είναι πυκνωτής σειράς με παράλληλο TCR, οπότε η ισοδύναμη αντίδραση
      σειράς της γραμμής μεταβάλλεται συνεχώς. Ο \textbf{SSSC} εγχέει τάση σε σειρά μέσω
      μετασχηματιστή σύζευξης σειράς· η εγχεόμενη τάση είναι ανεξάρτητη του ρεύματος της
      γραμμής, ενώ στον TCSC η τάση στα άκρα του πυκνωτή είναι ανάλογη του ρεύματος.
\item Ο \textbf{UPFC} αποτελείται από έναν παράλληλο και έναν σειριακό VSC με \emph{κοινό
      ζυγό ΣΡ} (DC link). Ο κοινός ζυγός επιτρέπει την ανταλλαγή \emph{ενεργού} ισχύος
      μεταξύ των δύο μετατροπέων, οπότε ο σειριακός μετατροπέας εγχέει τάση οποιουδήποτε
      μέτρου και γωνίας εντός των ονομαστικών του ορίων και ελέγχει ανεξάρτητα τις ροές
      ενεργού και αέργου ισχύος.
\end{itemize}

\begin{figure}[!htbp]
\centering
\includegraphics[width=\linewidth]{figures/facts-devices.svg}
\caption{Μονογραμμικά διαγράμματα των βασικών διατάξεων FACTS. Οι δύο πρώτες συνδέονται
παράλληλα, οι δύο επόμενες σε σειρά, ενώ ο UPFC συνδυάζει αμφότερους τους τρόπους μέσω
κοινού ζυγού ΣΡ.}
\label{fig:facts-devices}
\end{figure}

\subsection{Χαρακτηριστικές V--I των SVC και STATCOM}

Η διαφορά στην αρχή λειτουργίας αποτυπώνεται στις χαρακτηριστικές $V$--$I$ των δύο
διατάξεων παράλληλης σύνδεσης (Σχήμα~\ref{fig:facts-vi}).

Ο SVC συμπεριφέρεται ως μεταβλητή επιδεκτικότητα, οπότε $I = B\,V$ και $Q = B\,V^2$. Στο
όριο χωρητικής λειτουργίας το ρεύμα είναι ανάλογο της τάσης του ζυγού: στο 50\% της
ονομαστικής τάσης ο SVC αποδίδει το 50\% του ονομαστικού ρεύματος και το 25\% της ονομαστικής
αέργου ισχύος. Ο STATCOM, ως πηγή τάσης με έλεγχο ρεύματος, διατηρεί το ονομαστικό ρεύμα έως
πολύ χαμηλές τιμές τάσης, με περιορισμό μόνο από το μέγιστο επιτρεπόμενο ρεύμα των
ημιαγωγικών διακοπτών· η άεργος ισχύς του μειώνεται επομένως γραμμικά και όχι τετραγωνικά
με την τάση.

Κατά τη διάρκεια και αμέσως μετά από βραχυκύκλωμα, όταν η τάση είναι χαμηλή και η ανάγκη
στήριξης τάσης μέγιστη, ο STATCOM αποδίδει σημαντικά μεγαλύτερη άεργο ισχύ από SVC ίδιας
ονομαστικής ισχύος.

\begin{figure}[!htbp]
\centering
\includegraphics[width=\linewidth]{figures/facts-vi.svg}
\caption{Χαρακτηριστικά $V$--$I$ των δύο διατάξεων παράλληλης σύνδεσης. Η σκιασμένη περιοχή
είναι η περιοχή λειτουργίας και η έντονη γραμμή η χαρακτηριστική λειτουργίας με στατισμό. Τα
όρια του SVC είναι ευθείες σταθερής επιδεκτικότητας που διέρχονται από την αρχή, ενώ του
STATCOM είναι κατακόρυφες ευθείες σταθερού ρεύματος.}
\label{fig:facts-vi}
\end{figure}

Οι διατάξεις FACTS, μαζί με τα συστήματα συνεχούς ρεύματος υψηλής τάσης (High Voltage Direct
Current, HVDC), στα οποία η
ροή ισχύος καθορίζεται απευθείας από τον έλεγχο των μετατροπέων, αποτελούν τα κύρια μέσα
ρύθμισης τάσης και διαχείρισης συμφορήσεων στο σύστημα μεταφοράς. Τα αντίστοιχα μέσα στα
δίκτυα διανομής εξετάζονται στα Κεφάλαια 5 και 6.
